\chapter[Oszacowania Calderóna--Zygmunda i Schaudera]{Regularność $W^{2,p}$, całki osobliwe Calderóna--Zygmunda i oszacowania Schaudera}
\label{chap:pde16}

Rozdział~\ref{chap:pde15} dostarczył oszacowań drugich pochodnych w $L^2$.
Nie jest to koniec regularności. Dla potencjału Newtona z odpowiednimi danymi $f\in L^p$
uzyskamy oszacowania $D^2u\in L^p$ dla $1<p<\infty$,
a przy lokalnie Hölderowskich danych równania Poissona
wykażemy wewnętrzną regularność $C^{2,\alpha}$ rozwiązań klasy $C^2$.
Przejścia te nie wynikają z różniczkowania funkcji $f$: mechanizm
opiera się na strukturze całek osobliwych i na wzajemnym znoszeniu
się ich dodatnich i ujemnych części. Właśnie dlatego końcowe
przypadki $p=1$, $p=\infty$ oraz $\alpha=1$ są subtelne.

Przez $\Omega$ rozumiemy obszar w $\R^d$; $U\Subset V$ oznacza
$\overline U\subset V$ i zwartość $\overline U$. W tym rozdziale
$d\ge2$, chyba że jawnie powiemy inaczej. Wszystkie całki po
$\R^d$ i jego podzbiorach zapisujemy względem $\mathrm d\lambda_d$.
Przyjmujemy konwencję Fouriera z~\eqref{eq:pde8-convention}.

\section[Dlaczego pojawiają się całki osobliwe?]{Dlaczego $D^2(-\Delta)^{-1}$ nie jest operatorem splotu z funkcją całkowalną?}
Niech $\omega_{d-1}=\mathcal H^{d-1}(S^{d-1})$ oznacza pole
jednostkowej sfery, a $\Gamma$ rozwiązanie fundamentalne $-\Delta$:
$-\Delta\Gamma=\delta_0$. Z twierdzenia~\ref{thm:pde7-laplace} dla $d\ge3$ można wziąć
$\Gamma(x)=\bigl((d-2)\omega_{d-1}\bigr)^{-1}|x|^{2-d}$,
a dla $d=2$ funkcję $\Gamma(x)=-(2\pi)^{-1}\log|x|$.
Dla $x\ne0$ w każdym z obu przypadków
$\partial_i\Gamma(x)=-x_i/(\omega_{d-1}|x|^d)$.
Różniczkując względem $x_j$, otrzymujemy
\begin{equation}\label{eq:pde16-kernel}
 K_{ij}(x):=\partial_{ij}\Gamma(x)
 =\frac{d x_ix_j-\delta_{ij}|x|^2}
 {\omega_{d-1}|x|^{d+2}},\qquad 1\le i,j\le d,
\end{equation}
gdzie $\delta_{ij}$ oznacza deltę Kroneckera.
Jądro jest gładkie poza zerem i spełnia
$K_{ij}(tx)=t^{-d}K_{ij}(x)$ dla $t>0$.
Różniczkowanie tej tożsamości daje jednorodność gradientu
stopnia $-d-1$. Z ograniczoności obu funkcji na zwartej
sferze wynika $|K_{ij}(x)|\le C_d|x|^{-d}$ oraz
$|\nabla K_{ij}(x)|\le C_d|x|^{-d-1}$.
Dla $i\ne j$ odbicie względem $\theta_i=0$ zmienia znak
$\theta_i\theta_j$ i zachowuje miarę powierzchniową.
Dla $i=j$ całki z $\theta_i^2$ są równe przez symetrię
permutacji współrzędnych, a ich suma wynosi $\omega_{d-1}$.
Stąd
\[
 \int_{S^{d-1}}\theta_i\theta_j\,\mathrm d\mathcal H^{d-1}
 =\frac{\omega_{d-1}}d\delta_{ij},\qquad
 \int_{S^{d-1}}K_{ij}(\theta)\,\mathrm d\mathcal H^{d-1}(\theta)=0.
\]
Każde jądro jest niezerowe na pewnym otwartym podzbiorze
sfery (dla $i=j$ weźmy $\theta=e_i$, dla $i\ne j$
$\theta=(e_i+e_j)/\sqrt2$). Współrzędne biegunowe dają
\[
 \int_{B_1(0)}|K_{ij}(z)|\,\mathrm d\lambda_d(z)
 =\int_{S^{d-1}}|K_{ij}(\theta)|\,\mathrm d\mathcal H^{d-1}(\theta)
 \int_0^1 r^{-1}\,\mathrm d\lambda_1(r)=+\infty .
\]
Jeżeli $f$ jest ciągła i $f(x)\ne0$, to $|f(x-z)|$
jest oddzielone od zera w otoczeniu zera. Zwykła całka
Lebesgue'a z $K_{ij}(z)f(x-z)$ nie jest więc bezwzględnie
zbieżna, nawet dla takich $f\in C_c^\infty$.

\begin{definicja}[Główna wartość całki]
Niech $f\in C_c^1(\R^d)$, a $K$ będzie jądrem klasy $C^1$ poza
zerem, jednorodnym stopnia $-d$ i o zerowej średniej sferycznej.
Definiujemy
\[
 T_Kf(x)=\operatorname{p.v.}\!\int_{\R^d}K(x-y)f(y)\,\mathrm d\lambda_d(y)
 :=\lim_{\varepsilon\downarrow0}
 \int_{|x-y|>\varepsilon}K(x-y)f(y)\,\mathrm d\lambda_d(y).
\]
Dla $0<\varepsilon<1$ podstawiamy $z=x-y$.
Jednorodność i zerowa średnia dają zerową całkę z $K$
po każdym pierścieniu. Zatem po odjęciu $f(x)$
w pierścieniu $\varepsilon<|z|<1$ dostajemy w granicy
\[
 \begin{aligned}
 T_Kf(x)&=\int_{|z|<1}K(z)(f(x-z)-f(x))\,\mathrm d\lambda_d(z)\\
 &\quad+\int_{|z|\ge1}K(z)f(x-z)\,\mathrm d\lambda_d(z).
 \end{aligned}
\]
Obie całki są bezwzględnie zbieżne: w pierwszej
$|f(x-z)-f(x)|\le\|\nabla f\|_\infty |z|$,
więc całka radialna jest ograniczona przez stałą razy
$\int_0^1\mathrm d\lambda_1(r)$; druga ma moduł
nie większy niż $C\|f\|_1$. Uzasadnia to istnienie
granicy dla każdego $x$ oraz
$|T_Kf(x)|\le C(\|\nabla f\|_\infty+\|f\|_1)$.
\end{definicja}

Dystrybucję $\operatorname{p.v.}K$ określamy przez
\[
 \begin{aligned}
 \langle\operatorname{p.v.}K,\varphi\rangle
 &=\int_{|z|<1}K(z)(\varphi(z)-\varphi(0))\,\mathrm d\lambda_d(z)\\
 &\quad+\int_{|z|\ge1}K(z)\varphi(z)\,\mathrm d\lambda_d(z).
 \end{aligned}
\]
Dla testów o nośniku w ustalonym zwartym $L$ moduł
działania jest ograniczony przez
$C\|\nabla\varphi\|_\infty+C\lambda_d(L)\|\varphi\|_\infty$.
Liniowość i oszacowanie przez półnormy na każdym zwartym
nośniku dowodzą ciągłości dystrybucyjnej.

\begin{lemat}[Hessian rozwiązania fundamentalnego jako dystrybucja]
\label{lem:pde16-distr-kernel}
Dla każdego $1\le i,j\le d$ zachodzi w $\mathcal D'(\R^d)$
\begin{equation}\label{eq:pde16-distr-kernel}
 \partial_{ij}\Gamma
 =\operatorname{p.v.}K_{ij}-\frac{\delta_{ij}}d\delta_0.
\end{equation}
W szczególności dla $f\in C_c^\infty(\R^d)$ i $u=\Gamma*f$,
\begin{equation}\label{eq:pde16-repr}
 \partial_{ij}u=T_{ij}f-\frac{\delta_{ij}}d f,
 \qquad T_{ij}=T_{K_{ij}}.
\end{equation}
\end{lemat}
\begin{proof}
Niech $\varphi\in C_c^\infty(\R^d)$, $R>1$ tak duże,
aby nośnik $\varphi$ leżał w $B_R(0)$, oraz
$D_\varepsilon=\{x:\varepsilon<|x|<R\}$.
Z lokalnej całkowalności $\Gamma$ działanie
$\partial_{ij}\Gamma$ jest granicą całki
$\int_{D_\varepsilon}\Gamma\partial_{ij}\varphi\,\mathrm d\lambda_d$.
Całkując przez części najpierw względem $x_i$, potem
$x_j$, otrzymujemy
\[
 \begin{aligned}
 \int_{D_\varepsilon}\Gamma\partial_{ij}\varphi\,\mathrm d\lambda_d
 &=\int_{D_\varepsilon}K_{ij}\varphi\,\mathrm d\lambda_d\\
 &\quad+\int_{\partial D_\varepsilon}
 \Gamma n_i\partial_j\varphi\,\mathrm d\mathcal H^{d-1}\\
 &\quad-\int_{\partial D_\varepsilon}
 (\partial_i\Gamma)n_j\varphi\,\mathrm d\mathcal H^{d-1}.
 \end{aligned}
\]
Na zewnętrznej sferze oba wyrazy brzegowe znikają.
Na wewnętrznej sferze $x=\varepsilon\theta$ mamy
$n=-\theta$, $\partial_i\Gamma(x)=
-\theta_i/(\omega_{d-1}\varepsilon^{d-1})$,
a miara powierzchniowa skaluje się przez
$\varepsilon^{d-1}$. Ostatni wyraz wynosi więc
\[
 -\frac1{\omega_{d-1}}\int_{S^{d-1}}\theta_i\theta_j
 \varphi(\varepsilon\theta)\,\mathrm d\mathcal H^{d-1}(\theta)
 \longrightarrow-\frac{\delta_{ij}}d\varphi(0).
\]
Przedostatni wyraz jest ograniczony co do modułu przez
$C\|\partial_j\varphi\|_\infty\varepsilon|\log\varepsilon|$
dla $d=2$ i przez $C_d\|\partial_j\varphi\|_\infty\varepsilon$
dla $d\ge3$, więc znika w granicy.
Pierwsza całka dąży do działania $\operatorname{p.v.}K_{ij}$.
Daje to~\eqref{eq:pde16-distr-kernel}.
Z lematu~\ref{lem:pde7-convolution-derivatives} splot
$u=\Gamma*f$ jest gładki, $-\Delta u=f$ oraz
$\partial_{ij}u=(\partial_{ij}\Gamma)*f$.
Działanie $\operatorname{p.v.}K_{ij}$ na teście
$z\mapsto f(x-z)$ jest dokładnie $T_{ij}f(x)$,
a $\delta_0*f=f$. Otrzymujemy~\eqref{eq:pde16-repr}.
\end{proof}

\section[Oszacowanie Hilbertowskie i rozkład Calderóna--Zygmunda]{Oszacowanie $L^2$ i rozkład Calderóna--Zygmunda}
\begin{lemat}[Oszacowanie w $L^2$]
\label{lem:pde16-l2}
Dla $f\in C_c^\infty(\R^d)$ zachodzi
$\|T_{ij}f\|_2\le C_d\|f\|_2$. Operator $T_{ij}$
przedłuża się jednoznacznie do ograniczonego operatora na
$L^2(\R^d)$.
\end{lemat}
\begin{proof}
Najpierw wykażemy przynależność do $L^2$, bez jednostajnego
oszacowania normy przez $\|f\|_2$. Z definicji głównej
wartości $T_{ij}f$ jest ograniczona na całej przestrzeni.
Niech $\operatorname{supp}f\subset B_R(0)$, $R\ge1$.
Dla $|x|>2R$ mamy $|x-y|\ge|x|/2$ na nośniku $f$,
więc całka nie jest osobliwa i
$|T_{ij}f(x)|\le C_d|x|^{-d}\|f\|_1$.
Całka kwadratu tego ograniczenia poza $B_{2R}$ jest
skończona, gdyż $\int_{2R}^\infty r^{-d-1}
\,\mathrm d\lambda_1(r)<\infty$.
Na $B_{2R}$ wystarcza ograniczoność. Stąd $T_{ij}f\in L^2$,
a z~\eqref{eq:pde16-repr} także $H_{ij}:=\partial_{ij}u\in L^2$.

Z $-\Delta u=f$ i przemienności pochodnych mamy
$-\Delta H_{ij}=\partial_{ij}f$ dystrybucyjnie.
Obie strony są dystrybucjami temperowanymi; równość
w $\mathcal S'$ wynika z gęstości $C_c^\infty$ w $\mathcal S$
(twierdzenie~\ref{thm:pde8-density}). Twierdzenie~\ref{thm:pde8-dist-rules} daje
\[
 |\xi|^2\widehat H_{ij}=-\xi_i\xi_j\widehat f
 \qquad\text{w }\mathcal S'(\R^d).
\]
Obie strony są funkcjami lokalnie całkowalnymi, ponieważ
$\widehat H_{ij},\widehat f\in L^2$ z twierdzenia
Plancherela. Z jednoznaczności dystrybucji regularnych
(stwierdzenie~\ref{prop:pde7-injective}) mamy równość
prawie wszędzie. Dzieląc dla $\xi\ne0$, otrzymujemy
$\widehat H_{ij}(\xi)=-(\xi_i\xi_j/|\xi|^2)\widehat f(\xi)$
prawie wszędzie. Nie ma dodatkowego składnika skupionego
w zerze, bo $\widehat H_{ij}$ jest funkcją $L^2$.
Z~\eqref{eq:pde16-repr} wynika
$\widehat{T_{ij}f}=m_{ij}\widehat f$, gdzie
$m_{ij}(\xi)=\delta_{ij}/d-\xi_i\xi_j/|\xi|^2$
poza zerem, a wartość w zerze jest nieistotna.
Mamy $|m_{ij}|\le1+1/d\le2$. Z twierdzenia
~\ref{thm:pde8-plancherel} przy przyjętej konwencji
\[
 \|T_{ij}f\|_2=(2\pi)^{-d/2}\|m_{ij}\widehat f\|_2
 \le 2(2\pi)^{-d/2}\|\widehat f\|_2=2\|f\|_2.
\]
Argument działa również dla $d=2$, bez wymagania, aby
sam potencjał $u$ należał do $L^2$. Liniowość, gęstość
$C_c^\infty$ w $L^2$ i zupełność $L^2$ dają jednoznaczne
ciągłe przedłużenie (gęstość w $L^2$: twierdzenie
~\ref{thm:pde8-density}).
\end{proof}

\begin{lemat}[Rozkład Calderóna--Zygmunda]
\label{lem:pde16-cz-decomp}
Niech $f\in L^1(\R^d)\cap L^2(\R^d)$ i $s>0$. Istnieje
co najwyżej przeliczalna rodzina parami rozłącznych półotwartych
kostek dyadycznych $Q_k$ oraz rozkład $f=g+\sum_k b_k$
zbieżny w $L^1$ i $L^2$, taki że
\begin{enumerate}[label=(\roman*)]
\item $|g|\le 2^ds$ prawie wszędzie i $\|g\|_1\le\|f\|_1$;
\item $b_k=0$ prawie wszędzie poza $Q_k$,
 $\int_{Q_k} b_k\,\mathrm d\lambda_d=0$,
 $\sum_k\|b_k\|_1\le2\|f\|_1$;
\item $\sum_k|Q_k|\le s^{-1}\|f\|_1$ oraz
 $\|g\|_2^2\le 2^ds\|f\|_1$.
\end{enumerate}
\end{lemat}
\begin{proof}
Ustalamy siatkę kostek $2^m(k+[0,1)^d)$,
gdzie $m\in\mathbb Z$, $k\in\mathbb Z^d$.
Dwie kostki z tej siatki są rozłączne lub jedna
zawiera się w drugiej. Każda ma jedynego rodzica
o dwukrotnie większym boku. Dla $w\in L^1$
i $s>0$ nazywamy kostkę złą, jeśli średnia $|w|$
na niej przekracza $s$. Każda zła kostka zawiera
się w maksymalnej złej kostce: w łańcuchu rodziców
objętości rosną do nieskończoności, a średnie
nie przekraczają $\|w\|_1/|Q|$, więc tylko skończenie
wiele z nich może być złych. Maksymalne złe kostki
są parami rozłączne. Oznaczając przez $M_{\mathrm{dy}}w(x)$
supremum średnich $|w|$ po kostkach dyadycznych
zawierających $x$, dostajemy
\[
 \lambda_d(\{M_{\mathrm{dy}}w>s\})
 =\sum_{Q\ \mathrm{maks.}}|Q|
 \le s^{-1}\sum_{Q\ \mathrm{maks.}}\int_Q|w|\,\mathrm d\lambda_d
 \le s^{-1}\|w\|_1.
\]
Supremum jest mierzalne, gdyż siatka jest przeliczalna.

Uzasadnijmy także różniczkowanie całek po tych kostkach.
Niech $Q_n(x)$ oznacza jedyną kostkę o boku $2^{-n}$
zawierającą $x$ i $A_n w(x)$ średnią $w$ na $Q_n(x)$.
Dla $v\in C_c^\infty(\R^d)$ mamy $A_nv(x)\to v(x)$
w każdym punkcie przez ciągłość i średnicę kostki dążącą
do zera. Dla dowolnego $w\in L^1$ wybieramy
mierzalnego reprezentanta; zachodzi prawie wszędzie
\[
 \limsup_n |A_nw(x)-w(x)|
 \le M_{\mathrm{dy}}(w-v)(x)+|w(x)-v(x)|.
\]
Miara zbioru, gdzie lewa strona przekracza $\tau>0$,
jest więc ograniczona przez $4\|w-v\|_1/\tau$,
na mocy powyższej nierówności i nierówności Czebyszewa.
Z gęstości $C_c^\infty$ w $L^1$
(twierdzenie~\ref{thm:pde8-density}) miara ta wynosi zero.
Biorąc $\tau=1/l$ dla $l\in\N$, otrzymujemy
$A_nw\to w$ prawie wszędzie.

Wracamy do $w=f$ i oznaczamy maksymalne złe kostki
przez $Q_k$. Ich rodzice nie są źli, więc
\[
 s<\frac1{|Q_k|}\int_{Q_k}|f|\,\mathrm d\lambda_d
 \le 2^ds.
\]
Poza sumą tych kostek wszystkie średnie $|f|$ są
nie większe niż $s$; zastosowane do $|f|$ różniczkowanie
daje $|f|\le s$ prawie wszędzie. Kładziemy
$g=f$ poza kostkami oraz $g=f_{Q_k}$ na $Q_k$, gdzie
$f_{Q_k}=|Q_k|^{-1}\int_{Q_k}f\,\mathrm d\lambda_d$,
a $b_k=(f-f_{Q_k})\chi_{Q_k}$.
Z definicji średniej wynika zerowa całka $b_k$.
Ponadto $|f_{Q_k}|\le 2^ds$ i
\[
 \|g\|_1
 \le\int_{\R^d\setminus\bigcup_k Q_k}|f|\,\mathrm d\lambda_d
    +\sum_k\int_{Q_k}|f|\,\mathrm d\lambda_d
 =\|f\|_1,\qquad
 \|b_k\|_1\le2\int_{Q_k}|f|\,\mathrm d\lambda_d.
\]
Otrzymujemy $\|g\|_2^2\le2^ds\|f\|_1$ oraz
oszacowanie sumy norm $b_k$. Oszacowanie objętości
wynika z nierówności dla $M_{\mathrm{dy}}f$.
Każdy $b_k$ należy do $L^2$, a punktowo prawie wszędzie
$\sum_k b_k=f-g$, przy czym w każdym punkcie co najwyżej
jeden składnik jest niezerowy.
Dla $q=1,2$ z rozłączności i twierdzenia Tonellego
\[
 \sum_k\|b_k\|_q^q=\|f-g\|_q^q<\infty .
\]
Potęga $q$ normy ogona jest sumą odpowiednich potęg
norm, więc szeregi zbiegają w obu przestrzeniach.
Wybierając reprezentanta $b_k$ równego zero poza $Q_k$,
mamy jego nośnik topologiczny zawarty w $\overline{Q_k}$,
niekoniecznie w samej półotwartej kostce.
\end{proof}

\begin{twierdzenie}[Słaby typ $(1,1)$ operatorów $T_{ij}$]
\label{thm:pde16-weak11}
Istnieje $C_d>0$ taka, że dla każdego
$f\in L^1(\R^d)\cap L^2(\R^d)$ i $t>0$ zachodzi
\begin{equation}\label{eq:pde16-weak11}
 \lambda_d\bigl(\{x:|T_{ij}f(x)|>t\}\bigr)
 \le \frac{C_d}{t}\,\|f\|_{L^1(\R^d)}.
\end{equation}
\end{twierdzenie}
\begin{proof}
Piszemy $T=T_{ij}$, $K=K_{ij}$ i stosujemy
lemat~\ref{lem:pde16-cz-decomp} z $s=t$.
Niech $Q_k^*$ oznacza domkniętą kostkę o tym samym środku
$c_k$ co $Q_k$ i boku $3\sqrt d\,\ell_k$,
gdzie $\ell_k$ jest bokiem $Q_k$. Dla
$E=\bigcup_k Q_k^*$ mamy
$|E|\le(3\sqrt d)^d\sum_k|Q_k|\le C_d\|f\|_1/t$.
Dla części dobrej nierówność Czebyszewa i
lemat~\ref{lem:pde16-l2} dają
\[
 \lambda_d(\{|Tg|>t/2\})
 \le 4t^{-2}\|Tg\|_2^2
 \le C_dt^{-1}\|f\|_1.
\]

Uzasadnijmy reprezentację przedłużonego operatora $T$
dla $b_k$, które nie muszą być gładkie.
Wygładzamy $b_k$ standardowymi jądrami
$\rho_\varepsilon$ o nośnikach w $B_\varepsilon(0)$.
Funkcje $b_k*\rho_\varepsilon\in C_c^\infty$
zbiegają do $b_k$ w $L^1$ i $L^2$; wynika to
z twierdzenia~\ref{thm:pde7-approx-id} na ustalonym
zwartym otoczeniu $\overline{Q_k}$, gdyż wszystkie
nośniki są w nim zawarte dla małych $\varepsilon$.
Z ciągłości $T$ na $L^2$ i wyboru podciągu
$T(b_k*\rho_\varepsilon)\to Tb_k$ prawie wszędzie.
Dla ustalonego $x\notin Q_k^*$ odległość od
$\overline{Q_k}$ jest dodatnia; dla małych
$\varepsilon$ jądro $K(x-y)$ jest ograniczone
na wspólnym zwartym otoczeniu nośników.
Zbieżność w $L^1$ daje tam granicę zwykłych całek.
Stąd dla prawie każdego $x\notin Q_k^*$, dzięki zerowej
całce $b_k$,
\[
 Tb_k(x)=\int_{Q_k}[K(x-y)-K(x-c_k)]
 b_k(y)\,\mathrm d\lambda_d(y).
\]
Dla $y\in Q_k$ mamy $|y-c_k|\le\sqrt d\,\ell_k/2$,
a dla $x\notin Q_k^*$ zachodzi
$|x-c_k|\ge3\sqrt d\,\ell_k/2$.
Każdy punkt odcinka od $x-c_k$ do $x-y$ ma więc
moduł co najmniej $2|x-c_k|/3$. Stosując podstawowe
twierdzenie rachunku całkowego wzdłuż odcinka oraz
oszacowanie gradientu jądra, dostajemy
\[
 |Tb_k(x)|\le C_d\int_{Q_k}
 \frac{|y-c_k|\,|b_k(y)|}{|x-c_k|^{d+1}}
 \,\mathrm d\lambda_d(y)
 \le C_d\ell_k |x-c_k|^{-d-1}\|b_k\|_1.
\]
Całkując i powiększając zbiór całkowania do
$\{|x-c_k|\ge3\sqrt d\,\ell_k/2\}$, otrzymujemy
\[
 \int_{\R^d\setminus Q_k^*}|Tb_k|\,\mathrm d\lambda_d
 \le C_d\ell_k\|b_k\|_1
 \int_{3\sqrt d\,\ell_k/2}^\infty r^{-2}\,\mathrm d\lambda_1(r)
 \le C_d\|b_k\|_1.
\]
Dla skończonych sum nierówność trójkąta daje
$\int_{\R^d\setminus E}|T(\sum_{k\le N}b_k)|
\,\mathrm d\lambda_d\le C_d\sum_k\|b_k\|_1$.
Ponieważ $\sum_{k\le N}b_k\to f-g$ w $L^2$,
obrazy zbiegają w $L^2$ do $T(f-g)$.
Wybieramy podciąg zbieżny prawie wszędzie
i stosujemy lemat Fatou. Otrzymujemy
\[
 \int_{\R^d\setminus E}|T(f-g)|\,\mathrm d\lambda_d
 \le C_d\sum_k\|b_k\|_1\le C_d\|f\|_1.
\]
Z liniowości $Tf=Tg+T(f-g)$ prawie wszędzie.
Zbiór $\{|Tf|>t\}$ zawiera się z dokładnością
do zbioru miary zero w sumie $E$,
$\{|Tg|>t/2\}$ i
$\{x\notin E:|T(f-g)(x)|>t/2\}$.
Ostatnią miarę szacujemy przez $2/t$ razy
poprzednią całkę. To kończy dowód.
\end{proof}

\begin{lemat}[Interpolacja pomiędzy słabym $L^1$ i silnym $L^2$]
\label{lem:pde16-interp}
Niech $T$ będzie ograniczonym operatorem liniowym
na $L^2(\R^d)$, spełniającym oszacowanie słabego typu
$(1,1)$ na $L^1\cap L^2$ ze stałą $C_1$.
Dla każdego $1<p<2$ jego ograniczenie do
$L^1\cap L^2\cap L^p$ spełnia
$\|Tf\|_p\le C_p\|f\|_p$, gdzie $C_p$ zależy również
od stałych obu oszacowań. Ma ono jednoznaczne
ograniczone przedłużenie na $L^p$.
\end{lemat}
\begin{proof}
Dla $f\in L^1\cap L^2\cap L^p$ i $t>0$ kładziemy
$f_t^{>}=f\chi_{\{|f|>t\}}$ i
$f_t^{\le}=f\chi_{\{|f|\le t\}}$.
Obie funkcje należą do $L^1\cap L^2$.
Z liniowości i nierówności trójkąta wynika
$\{|Tf|>2t\}\subset\{|Tf_t^{>}|>t\}
\cup\{|Tf_t^{\le}|>t\}$ z dokładnością do zbioru
miary zero. Oznaczając $\|T\|_{L^2\to L^2}=C_2$,
otrzymujemy
\[
 \lambda_d(\{|Tf|>2t\})
 \le \frac{C_1}{t}\int_{\{|f|>t\}}|f|\,\mathrm d\lambda_d
 +\frac{C_2^2}{t^2}\int_{\{|f|\le t\}}|f|^2\,\mathrm d\lambda_d .
\]
Tożsamość
$|Tf(x)|^p=\int_0^{|Tf(x)|}p s^{p-1}\,\mathrm d\lambda_1(s)$
i twierdzenie Tonellego dają tożsamość warstwową
\[
 \|Tf\|_p^p=p\int_0^\infty s^{p-1}
 \lambda_d(\{|Tf|>s\})\,\mathrm d\lambda_1(s)
\]
również wtedy, gdy początkowo obie strony są nieskończone.
Po podstawieniu $s=2t$ i użyciu Tonellego otrzymujemy
\[
 \begin{aligned}
 \|Tf\|_p^p
 &\le p2^p C_1\int_{\{|f|>0\}} |f(x)|
       \int_0^{|f(x)|}t^{p-2}\,\mathrm d\lambda_1(t)
       \,\mathrm d\lambda_d(x)\\
 &\quad+p2^p C_2^2\int_{\{|f|>0\}} |f(x)|^2
       \int_{|f(x)|}^\infty t^{p-3}\,\mathrm d\lambda_1(t)
       \,\mathrm d\lambda_d(x)\\
 &=p2^p\left(\frac{C_1}{p-1}+\frac{C_2^2}{2-p}\right)
   \|f\|_p^p .
 \end{aligned}
\]
Punkty $f=0$ pomijamy, gdyż pierwotne całki po $t$
są w nich zerowe. Warunki $p>1$ i $p<2$ odpowiadają
zbieżności obu całek wewnętrznych.
Gęstość $C_c^\infty$ w $L^p$
(twierdzenie~\ref{thm:pde8-density}) i zupełność $L^p$
dają jednoznaczne przedłużenie.
Jest ono zgodne z wyjściowym operatorem $T$ na
$L^2\cap L^p$: dla $f$ z tej przestrzeni funkcje
$f_n=f\chi_{B_n(0)}\chi_{\{|f|\le n\}}$
należą do $L^1\cap L^2\cap L^p$ i zbiegają do $f$
w obu normach przez zbieżność zmajoryzowaną.
Ich obrazy zbiegają w $L^2$ do $Tf$ oraz w $L^p$
do wartości przedłużenia. Z obu zbieżności wybieramy
wspólny podciąg zbieżny prawie wszędzie, więc granice
są równe prawie wszędzie.
\end{proof}

\begin{twierdzenie}[Oszacowanie Calderóna--Zygmunda dla laplasjanu]
\label{thm:pde16-cz}
Niech $d\ge2$ i $1<p<\infty$. Dla
$f\in C_c^\infty(\R^d)$ oraz $u=\Gamma*f$ zachodzi
\begin{equation}\label{eq:pde16-cz}
 \|D^2u\|_{L^p(\R^d)}\le C_{d,p}\|f\|_{L^p(\R^d)}.
\end{equation}
Oszacowanie przedłuża operator $f\mapsto D^2(\Gamma*f)$
do $L^p$; nie oznacza ono, że samo $u$ musi należeć
do $L^p(\R^d)$.
\end{twierdzenie}
\begin{proof}
Z twierdzenia~\ref{thm:pde16-weak11},
lematu~\ref{lem:pde16-l2} i
lematu~\ref{lem:pde16-interp} wynika ograniczoność
$T_{ij}$ w $L^p$ dla $1<p\le2$.
Dla $p>2$ kładziemy $p'=p/(p-1)\in(1,2)$.
Mnożnik $m_{ij}$ wyprowadzony w dowodzie
lematu~\ref{lem:pde16-l2} jest rzeczywisty.
Z tożsamości Parsevala--Plancherela
~\ref{thm:pde8-plancherel} wynika więc
\[
 \int_{\R^d}(T_{ij}f)\overline h\,\mathrm d\lambda_d
 =\int_{\R^d}f\,\overline{T_{ij}h}\,\mathrm d\lambda_d
 \qquad(f,h\in L^2).
\]
Najpierw bierzemy $f,h\in C_c^\infty$.
Z nierówności Höldera i już udowodnionej ograniczoności
w $L^{p'}$ dostajemy
\[
 \left|\int_{\R^d}(T_{ij}f)\overline h\,\mathrm d\lambda_d\right|
 \le\|f\|_p\|T_{ij}h\|_{p'}
 \le C_{d,p'}\|f\|_p\|h\|_{p'}.
\]
Dla ustalonego testu $f$ wiemy już, że $T_{ij}f$ jest
ograniczona i zanika jak $O(|x|^{-d})$ w nieskończoności
(z dowodu lematu~\ref{lem:pde16-l2}).
Stąd należy do $L^p$ dla każdego $p>1$, gdyż
$\int_1^\infty r^{d-1-dp}\,\mathrm d\lambda_1(r)<\infty$.
Wolno zatem obliczyć jej normę przez dualność z $L^{p'}$.
Dokładniej, dla niezerowej funkcji $q=T_{ij}f\in L^p$
funkcja $h_0=|q|^{p-2}q/\|q\|_p^{p-1}$ ma normę
$L^{p'}$ równą jeden i
$\int q\overline{h_0}\,\mathrm d\lambda_d=\|q\|_p$.
Nierówność Höldera daje oszacowanie przeciwne dla
dowolnego $h$ o normie co najwyżej jeden.
Gęstość testów w $L^{p'}$ (twierdzenie
~\ref{thm:pde8-density}) i w razie potrzeby
normalizacja przybliżeń $h_0$ pozwalają ograniczyć
supremum do testów o normie nie większej niż jeden.
Otrzymujemy $\|T_{ij}f\|_p\le C_{d,p'}\|f\|_p$,
a następnie przedłużenie przez gęstość.
Także dla $p>2$ przedłużenie jest zgodne z operatorem
$L^2$ na $L^p\cap L^2$. Istotnie, obcinamy funkcję
jak w dowodzie lematu~\ref{lem:pde16-interp},
a każde obcięcie wygładzamy w obu normach
(twierdzenie~\ref{thm:pde7-approx-id} i wspólny
zwarty nośnik). Wybierając dla $n$ parametr wygładzania
tak, aby oba błędy były mniejsze niż $1/n$,
otrzymujemy ciąg testów zbieżny w obu normach.
Zbieżność obrazów i wspólny podciąg prawie wszędzie
dają zgodność przedłużeń.
Ze wzoru~\eqref{eq:pde16-repr} mamy
$\|\partial_{ij}u\|_p\le(C_{d,p}+1/d)\|f\|_p$.
Norma Frobeniusa Hessianu jest punktowo nie większa
niż suma modułów jego $d^2$ składowych, więc nierówność
trójkąta w $L^p$ daje~\eqref{eq:pde16-cz}.
Zupełność przestrzeni pól macierzowych $L^p$
uzasadnia ostatnie zdanie tezy.
\end{proof}

\begin{uwaga}[Dlaczego przypadki $1$ i $\infty$ są osobne?]
\label{rem:pde16-endpoints}
Silne oszacowanie $L^1\to L^1$ jest fałszywe nawet
dla gładkich danych. Istotnie, wybierzmy
$f\in C_c^\infty(B_1(0))$, $f\ge0$, $\int f\,\mathrm d\lambda_d=1$.
Dla $|x|>2$ mamy
\[
 |T_{12}f(x)-K_{12}(x)|
 \le C_d|x|^{-d-1}\int |y|f(y)\,\mathrm d\lambda_d(y)
 \le C_d|x|^{-d-1}.
\]
Wynika to z oszacowania gradientu jądra na odcinku
od $x$ do $x-y$, którego punkty mają moduł co najmniej
$|x|/2$. Na pewnym otwartym fragmencie sfery wokół
$(e_1+e_2)/\sqrt2$ mamy $K_{12}(\theta)\ge c_d>0$.
Dzięki jednorodności, dla odpowiednio dużego $r$ i takich
$\theta$, $T_{12}f(r\theta)\ge(c_d/2)r^{-d}$.
Całka po tym stożku jest rozbieżna, ponieważ
$\int_R^\infty r^{-1}\,\mathrm d\lambda_1(r)=\infty$.
Ponieważ $\partial_{12}(\Gamma*f)=T_{12}f$,
nie zachodzi też silne oszacowanie Hessianu w $L^1$.
Oszacowanie~\eqref{eq:pde16-weak11} jest słabszym
oszacowaniem końcowym.

Dla końca $p=\infty$ w wymiarze $d=2$ weźmy
$f=\chi_{(0,1)^2}$ i $u=\Gamma*f$.
Poza domkniętym kwadratem wolno różniczkować pod całką
(lokalnie odległość od nośnika jest dodatnia).
Dla $x=(-t,-t)$, $0<t<1/8$, otrzymujemy
\[
 \partial_{12}u(x)=\frac1\pi
 \int_{(0,1)^2}
 \frac{(t+y_1)(t+y_2)}
 {((t+y_1)^2+(t+y_2)^2)^2}\,\mathrm d\lambda_2(y).
\]
Całkowana funkcja jest nieujemna. Na podsektorze
$\pi/6<\theta<\pi/3$, $2t<r<1/2$, gdzie
$y=(r\cos\theta,r\sin\theta)$, obie współrzędne $y$
są co najmniej $r/2$, a $|x-y|\le r+\sqrt2t\le2r$.
Stąd całkowana funkcja jest nie mniejsza niż $c r^{-2}$,
z $c>0$ niezależnym od $t$. Zatem
$\partial_{12}u(-t,-t)\ge c'\log(1/(4t))\to+\infty$.
Pochodna jest ciągła w otwartym dopełnieniu kwadratu,
więc każda dostatecznie duża wartość daje także
przekroczenie zadanego progu na otwartym zbiorze
dodatniej miary. Jest to brak ograniczoności istotnej,
a nie jedynie rozbieżność na ciągu punktów.
Przykład można przenieść na gładkie dane: wygładzenia
$f_\varepsilon=f*\rho_\varepsilon$ mają normę
$L^\infty$ nie większą niż $1$, a w każdym ustalonym
punkcie poza domkniętym kwadratem
$\partial_{12}(\Gamma*f_\varepsilon)$ zbiega
do powyższej całki dzięki zbieżności w $L^1$
i temu, że dla dostatecznie małych $\varepsilon$
wszystkie nośniki leżą w ustalonym zwartym otoczeniu
kwadratu, oddalonym od rozważanego punktu.
Jednostajne oszacowanie $L^\infty$ dla gładkich danych
prowadziłoby zatem do sprzeczności.
\end{uwaga}

\section[Lokalna regularność Sobolewa i zamrażanie współczynników]{Lokalna regularność $W^{2,p}$ i zamrażanie współczynników}
W tej sekcji kule $B_r,B_R$ o promieniach $0<r<R$
są współśrodkowe. Normę macierzy drugich pochodnych
liczymy za pomocą normy Frobeniusa.

\begin{lemat}[Lokalizacja oszacowania dla laplasjanu]
\label{lem:pde16-localization}
Niech $1<p<\infty$, $B_r\Subset B_R\subset\R^d$ i
$u\in W^{2,p}(B_R)$. Wówczas
\[
 \|D^2u\|_{L^p(B_r)}\le C_{d,p}\bigl(
 \|\Delta u\|_{L^p(B_R)}+
 (R-r)^{-2}\|u\|_{L^p(B_R)}+
 (R-r)^{-1}\|\nabla u\|_{L^p(B_R)}\bigr).
\]
\end{lemat}
\begin{proof}
Najpierw wyprowadzimy z twierdzenia~\ref{thm:pde16-cz}
oszacowanie
\[
 \|D^2v\|_{L^p(\R^d)}\le C_{d,p}\|\Delta v\|_{L^p(\R^d)}
 \quad(v\in W^{2,p}(\R^d),\ \operatorname{supp}v
 \text{ zwarty}).
\]
Dla $v\in C_c^\infty(\R^d)$ z lematu
~\ref{lem:pde7-convolution-derivatives} i tożsamości
$-\Delta\Gamma=\delta_0$ mamy
$\Gamma*(-\Delta v)=(-\Delta\Gamma)*v=v$.
Stosując~\eqref{eq:pde16-cz} do $f=-\Delta v$,
otrzymujemy żądane oszacowanie dla testów.
Dla $v$ o zwartej podporze wygładzenia
$v*\rho_\varepsilon$ są testami o wspólnym zwartym
nośniku i zbiegają do $v$ w $W^{2,p}(\R^d)$
na mocy twierdzenia~\ref{thm:pde9-mollify}
(zbieżność lokalna wystarcza dzięki wspólnemu nośnikowi).
W obu normach przechodzimy do granicy.

Niech $\sigma=(R-r)/4$ i $\rho_\sigma$ będzie
nieujemnym standardowym jądrem wygładzającym
o całce $1$ i nośniku w $B_\sigma(0)$.
Kładziemy $\eta=\chi_{B_{(r+R)/2}}*\rho_\sigma$.
Wówczas $0\le\eta\le1$, $\eta=1$ w otoczeniu
$\overline{B_r}$, a nośnik $\eta$ mieści się
w $\overline{B_{(r+R)/2+\sigma}}\subset B_R$.
Dla $|\beta|=1,2$ z różniczkowania splotu mamy
$\|\partial^\beta\eta\|_\infty\le
\|\partial^\beta\rho_\sigma\|_1
=\sigma^{-|\beta|}\|\partial^\beta\rho_1\|_1$.
Zatem $|D^k\eta|\le C_d(R-r)^{-k}$ dla $k=1,2$.
Z lematu~\ref{lem:pde9-compact-product} przedłużenie
$v$ funkcji $\eta u$ przez zero należy do
$W^{2,p}(\R^d)$ i ma zwarty nośnik. Reguła Leibniza
daje w $B_R$
\[
 \Delta v=\eta\Delta u+2\nabla\eta\cdot\nabla u+u\Delta\eta,
\]
a poza $B_R$ wszystkie składniki są zerowe
(po przedłużeniu przez zero).
Ponieważ $D^2v=D^2u$ na $B_r$, poprzednie oszacowanie,
nierówność trójkąta i ograniczenia pochodnych $\eta$
dają tezę.
\end{proof}

\begin{lemat}[Interpolacja pierwszych pochodnych]
\label{lem:pde16-gradinterp}
Niech $1<p<\infty$, $B_r\Subset B_R$, $u\in W^{2,p}(B_R)$.
Dla każdego $0<h<R-r$ zachodzi
\[
 \|\nabla u\|_{L^p(B_r)}
 \le \frac{2d}{h}\|u\|_{L^p(B_R)}
 +\frac{dh}{2}\|D^2u\|_{L^p(B_R)}.
\]
W szczególności dla każdego $\varepsilon>0$ mamy
\[
 \|\nabla u\|_{L^p(B_r)}\le
 \varepsilon\|D^2u\|_{L^p(B_R)}
 +C_{\varepsilon,r,R}\|u\|_{L^p(B_R)}.
\]
\end{lemat}
\begin{proof}
Dla funkcji gładkiej i $x\in B_r$ cały odcinek
$x+[0,h]e_j$ mieści się w $B_R$.
Dwukrotne użycie podstawowego twierdzenia rachunku
całkowego i zamiana kolejności całkowania dają
\[
 \begin{aligned}
 u(x+he_j)-u(x)
 &=\int_0^h\partial_ju(x+se_j)\,\mathrm d\lambda_1(s)\\
 &=h\partial_ju(x)+\int_0^h(h-t)
       \partial_{jj}u(x+te_j)\,\mathrm d\lambda_1(t).
 \end{aligned}
\]
Po podzieleniu przez $h$ i przeniesieniu drugiego
składnika otrzymujemy wzór na $\partial_ju(x)$.
Nierówność Minkowskiego dla całki i translacja zmiennych
dają
\[
 \|\partial_ju\|_{L^p(B_r)}
 \le \frac2h\|u\|_{L^p(B_{r+h})}
 +\frac1h\int_0^h(h-t)\,\mathrm d\lambda_1(t)
       \|\partial_{jj}u\|_{L^p(B_{r+h})}.
\]
Ponieważ całka z wagi wynosi $h^2/2$, po zsumowaniu
po $j$ mamy pierwsze oszacowanie, nawet z normami
na $B_{r+h}$ po prawej stronie. Dla dowolnego
$u\in W^{2,p}(B_R)$ wybieramy $R'$ z $r+h<R'<R$.
Z twierdzenia~\ref{thm:pde9-mollify} wygładzenia
$u_\varepsilon$ są gładkie na $B_{R'}$ dla małych
parametrów i zbiegają tam w $W^{2,p}$.
Stosujemy oszacowanie z prawą stroną na $B_{r+h}$
do wygładzeń i przechodzimy w normach do granicy.
Powiększenie prawej strony do $B_R$ daje żądany wzór.
Wybór $h=\min\{(R-r)/2,2\varepsilon/d\}$
daje drugie oszacowanie z $C_{\varepsilon,r,R}=2d/h$.
\end{proof}

\begin{twierdzenie}[Wewnętrzne oszacowanie dla współczynników ciągłych]
\label{thm:pde16-variable}
Niech $\Omega\subset\R^d$ będzie obszarem, $1<p<\infty$,
$a_{ij}\in C(\Omega)$ i niech symetryczna macierz
$A(x)=(a_{ij}(x))$ będzie lokalnie jednostajnie eliptyczna:
na każdym zwartym podzbiorze $\Omega$ istnieją
$0<\lambda\le\Lambda$ takie, że
$\lambda|\xi|^2\le\xi^{\mathsf T}A(x)\xi\le\Lambda|\xi|^2$
dla wszystkich $\xi\in\R^d$.
Piszemy $A:D^2u=\sum_{i,j=1}^d a_{ij}\partial_{ij}u$.
Jeżeli $u\in W^{2,p}_{\mathrm{loc}}(\Omega)$ i
$A:D^2u=f\in L^p_{\mathrm{loc}}(\Omega)$ prawie wszędzie,
to dla niepustych otwartych zbiorów $U\Subset V\Subset\Omega$
\begin{equation}\label{eq:pde16-variable}
 \|u\|_{W^{2,p}(U)}\le C\bigl(
 \|f\|_{L^p(V)}+\|u\|_{L^p(V)}\bigr).
\end{equation}
Stała zależy od $d,p,U,V$, eliptyczności i modułu ciągłości $A$
na $V$. Twierdzenie jest tu \emph{oszacowaniem a priori}: samo
w sobie nie konstruuje rozwiązania $W^{2,p}$.
\end{twierdzenie}
\begin{proof}
\emph{Krok 1: operator o stałych współczynnikach.}
Niech $A_0=A(x_0)$, $S=A_0^{1/2}$ i
$w(z)=v(x_0+Sz)$ dla $v\in C_c^\infty(\R^d)$.
Symetria i dodatnia określoność dają
$\|S\|_{\mathrm{op}}\le\sqrt\Lambda$ oraz
$\|S^{-1}\|_{\mathrm{op}}\le\lambda^{-1/2}$.
Klasyczna reguła łańcuchowa daje
$D_z^2w=S^{\mathsf T}(D_x^2v)(x_0+Sz)S$,
więc $\Delta_z w=(A_0:D_x^2v)(x_0+Sz)$.
Z oszacowania globalnego udowodnionego w lemacie
~\ref{lem:pde16-localization} i zmiany zmiennych
$x=x_0+Sz$ mamy
\[
 \begin{aligned}
 \|D_x^2v\|_p
 &\le\lambda^{-1}|\det S|^{1/p}\|D_z^2w\|_p\\
 &\le\lambda^{-1}C_{d,p}|\det S|^{1/p}\|\Delta_z w\|_p\\
 &=\lambda^{-1}C_{d,p}\|A_0:D_x^2v\|_p.
 \end{aligned}
\]
Wyznacznik się skraca. Przez wygładzanie jak
w dowodzie lematu~\ref{lem:pde16-localization}
oszacowanie obowiązuje także dla $v\in W^{2,p}(\R^d)$
o zwartym nośniku. Oznaczmy jego stałą przez $C_0$.

\emph{Krok 2: zamrażanie i lokalizacja.}
Wybieramy $L>0$ takie, że dla każdego
$x_0\in\overline U$ mamy $\overline{B_L(x_0)}\subset V$,
a ponadto
$\sup_{B_L(x_0)}|A(x)-A(x_0)|_{\mathrm F}\le1/(2C_0)$.
Taki wspólny promień istnieje z dodatniej odległości
$\overline U$ od dopełnienia $V$ i jednostajnej
ciągłości $A$ na $\overline V$.
Dla $v\in W^{2,p}(\R^d)$ o zwartym nośniku
w $B_L(x_0)$, z nierówności Cauchy'ego--Schwarza
dla iloczynu Frobeniusa, dostajemy
\[
 \begin{aligned}
 \|D^2v\|_p
 &\le C_0\|A:D^2v\|_{p,B_L}
      +C_0\|(A_0-A):D^2v\|_{p,B_L}\\
 &\le C_0\|A:D^2v\|_{p,B_L}+\tfrac12\|D^2v\|_p.
 \end{aligned}
\]
Po pochłonięciu mamy $\|D^2v\|_p\le
2C_0\|A:D^2v\|_{p,B_L}$.
Ustalmy środek $x_0$ i pomijajmy go w oznaczeniach kul.
Dla $0<s<t<L$ używamy odcięcia równego $1$ na $B_s$
i o zwartym nośniku w $B_t$, skonstruowanego w lemacie
~\ref{lem:pde16-localization}. Dla $v=\eta u$
przedłużonego przez zero reguła Leibniza daje,
dzięki symetrii $A$,
\[
 A:D^2v=\eta f+2(A\nabla\eta)\cdot\nabla u+u\,A:D^2\eta.
\]
Stąd, używając $\|A\|_{\mathrm{op}}\le\Lambda$
i $|A|_{\mathrm F}\le\sqrt d\,\Lambda$, otrzymujemy
\[
 \|D^2u\|_{p,B_s}\le C_1\bigl(
 \|f\|_{p,B_t}+(t-s)^{-1}\|\nabla u\|_{p,B_t}
 +(t-s)^{-2}\|u\|_{p,B_t}\bigr),
\]
gdzie $C_1$ jest niezależna od $s,t,x_0,u,f$.

\emph{Krok 3: pochłanianie na zagnieżdżonych kulach.}
Niech $H(s)=\|D^2u\|_{p,B_s}$,
$F=\|f\|_{p,B_L}$ i $Y=\|u\|_{p,B_L}$.
Dla $0<s<t<L$ stosujemy poprzednie oszacowanie
z promieniem pośrednim $m=(s+t)/2$.
Z lematu~\ref{lem:pde16-gradinterp}, dla
$h=\kappa(t-s)$, gdzie $0<\kappa<1/2$, mamy
\[
 \|\nabla u\|_{p,B_m}
 \le \frac{2d}{\kappa(t-s)}Y
       +\frac{d\kappa(t-s)}2 H(t).
\]
W konsekwencji
\[
 H(s)\le C_1 d\kappa H(t)+C_1 F
       +C_1(4+4d/\kappa)(t-s)^{-2}Y .
\]
Ustalamy $\kappa=\min\{1/4,1/(8dC_1)\}$,
tak że współczynnik przy $H(t)$ jest nie większy
niż $\theta=1/8$. Zatem z pewną stałą $C_2$,
\[
 H(s)\le\theta H(t)+C_2\bigl(F+(t-s)^{-2}Y\bigr).
\]
Bierzemy $s_n=L(1-2^{-n-1})$, $n\ge0$,
więc $s_0=L/2$ i $s_{n+1}-s_n=L2^{-n-2}$.
Iterując nierówność od $n=0$ do $N-1$, dostajemy
\[
 H(L/2)\le\theta^N H(s_N)
 +C_2 F\sum_{n=0}^{N-1}\theta^n
 +16C_2 L^{-2}Y\sum_{n=0}^{N-1}(4\theta)^n.
\]
Ponieważ $u\in W^{2,p}(B_L)$, mamy $H(s_N)\le H(L)<\infty$,
więc pierwszy składnik dąży do zera.
Oba szeregi geometryczne są zbieżne, gdyż $4\theta=1/2$.
Otrzymujemy
$H(L/2)\le C_3(F+L^{-2}Y)$.
Stosując lemat~\ref{lem:pde16-gradinterp}
do $B_{L/4}\Subset B_{L/2}$ z $h=L/8$, dostajemy
także $\|\nabla u\|_{p,B_{L/4}}\le
C_d(L^{-1}Y+L H(L/2))$.
Łącznie
\[
 \|u\|_{W^{2,p}(B_{L/4})}
 \le C_{d,p,\lambda,\Lambda,L}(F+Y).
\]

\emph{Krok 4: skończone pokrycie.}
Kule $B_{L/4}(x_0)$, $x_0\in\overline U$,
pokrywają zwarty zbiór $\overline U$; wybieramy
skończone podpokrycie. Na każdej kuli prawe strony
powyższych oszacowań są ograniczone przez wspólną
stałą razy $\|f\|_{p,V}+\|u\|_{p,V}$.
Całka po $U$ modułu do potęgi $p$ każdej pochodnej
rzędu co najwyżej dwa jest nie większa od sumy
analogicznych całek po kulach pokrycia.
Po zsumowaniu i wyciągnięciu pierwiastka stopnia $p$
otrzymujemy~\eqref{eq:pde16-variable}.
Pokrycie oraz $L$ wybrano niezależnie od $u,f$,
co daje deklarowaną zależność stałej.
\end{proof}

\begin{uwaga}[Postać dywergencyjna a niedywergencyjna]
Niech $A\in W^{1,\infty}_{\mathrm{loc}}$ będzie lokalnie
jednostajnie eliptyczna, $f\in L^2_{\mathrm{loc}}$ i
$u\in H^1_{\mathrm{loc}}$ niech spełnia
$-\diver(A\nabla u)=f$ słabo. Z twierdzenia
~\ref{thm:pde15-h2} mamy $u\in H^2_{\mathrm{loc}}$.
Reguła iloczynu dla mnożnika $W^{1,\infty}$,
uzasadniona w dowodzie twierdzenia
~\ref{thm:pde15-bootstrap}, daje wtedy prawie wszędzie
\[
 -\sum_{i,j=1}^d a_{ij}\partial_{ij}u
 =f+\sum_{i,j=1}^d(\partial_i a_{ij})\partial_j u .
\]
Obie strony należą lokalnie do $L^2$. Dla wykładnika
$p>2$ sama informacja $f\in L^p_{\mathrm{loc}}$
nie uzasadnia przynależności całej prawej strony
do $L^p_{\mathrm{loc}}$ na podstawie znanej jedynie
lokalnej regularności $H^2$. Ponadto nawet gdy
prawa strona należy do $L^p_{\mathrm{loc}}$, twierdzenie
~\ref{thm:pde16-variable} jest oszacowaniem a priori
dla funkcji już należących do $W^{2,p}_{\mathrm{loc}}$.
Do wykazania nowej regularności potrzebny jest odrębny
argument; nie można zakładać jej w dowodzie samej siebie.
Jeżeli $A$ nie jest symetryczna, w członie drugiego
rzędu można zastąpić ją przez $(A+A^{\mathsf T})/2$,
gdyż Hessian jest symetryczny prawie wszędzie.
\end{uwaga}

\section{Przestrzenie Höldera i regularność w przestrzeniach Höldera}
\begin{definicja}[Norma Höldera]
\label{def:pde16-holder}
Dla $0<\alpha<1$ i niepustego zbioru $U\subset\R^d$
definiujemy
\[
 [f]_{C^{0,\alpha}(U)}=
 \sup_{\substack{x,y\in U\\x\ne y}}
 \frac{|f(x)-f(y)|}{|x-y|^\alpha}.
\]
Supremum pustej rodziny przyjmujemy za zero.
Jeżeli $U$ jest ograniczonym zbiorem otwartym i
$k\in\N_0$, przestrzeń $C^{k,\alpha}(\overline U)$
składa się z funkcji klasy $C^k(U)$, których pochodne
do rzędu $k$ przedłużają się ciągle na $\overline U$,
a pochodne rzędu $k$ mają skończoną półnormę Höldera.
Normą jest
\[
 \|f\|_{C^{k,\alpha}(\overline U)}
 =\sum_{|\beta|\le k}\|\partial^\beta f\|_{\infty,U}
 +\sum_{|\beta|=k}[\partial^\beta f]_{C^{0,\alpha}(U)}.
\]
Półnormy i normy supremum można równoważnie liczyć
na domknięciu, przez gęstość $U$ w $\overline U$
i ciągłość przedłużeń. Na $\R^d$ przez
$C^{0,\alpha}(\R^d)$ rozumiemy przestrzeń funkcji
ograniczonych o skończonej globalnej półnormie Höldera,
z normą $\|f\|_\infty+[f]_\alpha$; indeks $c$
oznacza zwarty nośnik. Dla pól o wartościach w skończenie
wymiarowej przestrzeni stosujemy analogiczne definicje
z normą tej przestrzeni.
\end{definicja}

Będziemy używali nierówności
\[
 [ab]_\alpha\le\|a\|_\infty[b]_\alpha+
                    \|b\|_\infty[a]_\alpha,
\]
wynikającej z rozpisania
$a(x)b(x)-a(y)b(y)=a(x)(b(x)-b(y))+
b(y)(a(x)-a(y))$. W szczególności pełna norma
iloczynu jest nie większa od iloczynu pełnych norm.
Dla funkcji wektorowych i macierzowych analogiczne
oszacowania otrzymujemy składowo.

\begin{lemat}[Potencjał Newtona: oszacowanie Höldera]
\label{lem:pde16-newton-holder}
Niech $0<\alpha<1$, $f\in C_c^{0,\alpha}(\R^d)$,
$\operatorname{supp}f\subset B_R(0)$ i $u=\Gamma*f$.
Wtedy $u\in C^{2,\alpha}_{\mathrm{loc}}(\R^d)$,
$-\Delta u=f$ klasycznie i dla każdego $M>0$
\[
 \|u\|_{C^{2,\alpha}(\overline{B_M(0)})}
 \le C(d,\alpha,R,M)\|f\|_{C^{0,\alpha}(\R^d)}.
\]
\end{lemat}
\begin{proof}
\emph{Krok 1: przypadek gładkich danych i norma supremum.}
Zaczynamy od $f\in C_c^\infty$. Piszemy $K=K_{ij}$.
Równość~\eqref{eq:pde16-repr} oraz zerowa średnia jądra
po sferach dają
\begin{equation}\label{eq:pde16-cancel}
 \begin{aligned}
 \partial_{ij}u(x)&=-\frac{\delta_{ij}}d f(x)\\
 &\quad+\int_{|z|<1}K(z)[f(x-z)-f(x)]\,\mathrm d\lambda_d(z)\\
 &\quad+\int_{|z|\ge1}K(z)f(x-z)\,\mathrm d\lambda_d(z).
 \end{aligned}
\end{equation}
Pierwsza całka jest bezwzględnie zbieżna i ograniczona
co do modułu przez
$C_d[f]_\alpha\int_0^1 r^{\alpha-1}\,\mathrm d\lambda_1(r)
=C_d[f]_\alpha/\alpha$.
Druga ma moduł nie większy niż $C_d\|f\|_1$,
gdyż $|K(z)|\le C_d$ dla $|z|\ge1$.
Zatem dla wszystkich $x$ mamy
$|\partial_{ij}u(x)|\le C_{d,\alpha,R}
\|f\|_{C^{0,\alpha}(\R^d)}$.

\emph{Krok 2: półnorma Höldera.}
Niech $x,x'\in B_M(0)$ i $h=|x-x'|\in(0,1/4)$.
Wybieramy $L=R+M+2$ oraz $D=B_L(x)$, które
zawiera nośnik $f$. Niech $T=T_{ij}$.
Ze znoszenia całek jądra po pierścieniach mamy
\[
 \begin{aligned}
 Tf(x)&=\int_D K(x-y)(f(y)-f(x))\,\mathrm d\lambda_d(y),\\
 Tf(x')&=\int_D K(x'-y)(f(y)-f(x'))\,\mathrm d\lambda_d(y)
           +f(x')J,
 \end{aligned}
\]
gdzie $J=\operatorname{p.v.}\int_D K(x'-y)\,\mathrm d\lambda_d(y)$.
Obie całki z różnicami funkcji są bezwzględnie
zbieżne. Ponieważ
$B_{L-h}(x')\subset D\subset B_{L+h}(x')$,
wartość główna całki po pierwszej z tych kul jest
zerowa, a resztę ograniczamy przez
\[
 |J|\le C_d\int_{L-h}^{L+h}r^{-1}\,\mathrm d\lambda_1(r)
 \le C_d h/L .
\]
Rozdzielamy $D$ na $N=B_{2h}(x)$ i $E=D\setminus N$.
Całka z $K(x-y)(f(y)-f(x))$ po $N$ ma moduł
nie większy niż $C_{d,\alpha}[f]_\alpha h^\alpha$.
To samo zachodzi dla całki z $K(x'-y)(f(y)-f(x'))$
po $N$, gdyż $N\subset B_{3h}(x')$.

Na $E$ rozpisujemy różnicę całkowanych funkcji
jako sumę
\[
 [K(x-y)-K(x'-y)](f(y)-f(x))
       +K(x'-y)(f(x')-f(x)).
\]
Dla $y\in E$ odcinek między $x-y$ i $x'-y$
jest oddalony od zera o co najmniej $|x-y|/2$.
Oszacowanie gradientu jądra i twierdzenie
o wartości średniej dają
$|K(x-y)-K(x'-y)|\le C_d h|x-y|^{-d-1}$.
Pierwszy składnik po scałkowaniu ma zatem moduł
nie większy niż
\[
 C_d h[f]_\alpha\int_{2h}^{L}r^{\alpha-2}\,
          \mathrm d\lambda_1(r)
 \le \frac{C_d}{1-\alpha}[f]_\alpha h^\alpha .
\]
Dla drugiego składnika potrzebujemy oszacowania
całki z jądra, a nie całki z jego modułu na całym $E$.
Zachodzą inkluzje (z dokładnością do brzegów)
\[
 \{y:3h<|y-x'|<L-h\}\subset E
 \subset\{y:h<|y-x'|<L+h\}.
\]
Całka po wewnętrznym pierścieniu jest zerowa.
Pozostałe części leżą w pierścieniach o promieniach
$h,3h$ oraz $L-h,L+h$. Wobec tego
\[
 \left|\int_E K(x'-y)\,\mathrm d\lambda_d(y)\right|
 \le C_d\left(\log3+\log\frac{L+h}{L-h}\right)\le C_d'.
\]
Mnożenie przez $|f(x')-f(x)|\le[f]_\alpha h^\alpha$
daje wymagane oszacowanie. Wreszcie wyraz
$f(x')J$ ma moduł nie większy niż
$C_d\|f\|_\infty h/L\le C_d\|f\|_\infty h^\alpha$.
Uwzględniając składnik $-\delta_{ij}f/d$
z~\eqref{eq:pde16-cancel}, otrzymujemy
\[
 |\partial_{ij}u(x)-\partial_{ij}u(x')|
 \le C_{d,\alpha}\|f\|_{C^{0,\alpha}} h^\alpha .
\]
Dla $h\ge1/4$ dzielimy dwukrotne oszacowanie
supremum przez $h^\alpha\ge4^{-\alpha}$.
To kończy kontrolę drugich pochodnych.

\emph{Krok 3: niższe pochodne i przejście do granicy.}
Funkcje $\Gamma$ oraz $\partial_i\Gamma$ są lokalnie
całkowalne. Dla pochodnej pierwszego rzędu wynika
to z oszacowania $|\partial_i\Gamma(z)|\le C_d|z|^{1-d}$.
Pochodna dystrybucyjna $\Gamma$ jest reprezentowana
przez tę pochodną klasyczną poza zerem: całkowanie
przez części poza $B_\varepsilon(0)$ daje wyraz brzegowy
rzędu $\varepsilon$ dla $d\ge3$ lub
$\varepsilon|\log\varepsilon|$ dla $d=2$, który zanika.
Zatem dla gładkiego $f$,
$\partial_i u=(\partial_i\Gamma)*f$.
Dla $x\in B_M(0)$ i $y\in B_R(0)$ mamy
$|x-y|<M+R$. Wobec tego
\[
 \|u\|_{\infty,B_M}+\sum_i\|\partial_i u\|_{\infty,B_M}
 \le\left(\int_{B_{M+R}(0)}
 (|\Gamma|+\sum_i|\partial_i\Gamma|)\,\mathrm d\lambda_d\right)
 \|f\|_\infty .
\]

Teraz niech $f\in C_c^{0,\alpha}$.
Bierzemy $f_n=f*\rho_{1/n}$ z nieujemnym jądrem
o całce jeden. Wszystkie nośniki leżą w $B_{R+2}(0)$,
$\|f_n\|_\infty\le\|f\|_\infty$,
$[f_n]_\alpha\le[f]_\alpha$ (przez całkowanie
różnicy wartości $f$), a
$\|f_n-f\|_\infty\le n^{-\alpha}[f]_\alpha$.
Dla $u_n=\Gamma*f_n$ powyższe oszacowanie lokalnych
całek daje jednostajną zbieżność $u_n$ do $u$
oraz $\partial_i u_n$ do $(\partial_i\Gamma)*f$
na każdej zwartej kuli. Drugie pochodne $u_n$
są jednostajnie ograniczone i mają jednostajnie
ograniczone półnormy Höldera na każdej takiej kuli.
Z twierdzenia Arzelà--Ascoliego i wyboru przekątniowego
na kulach o całkowitych promieniach otrzymujemy
podciąg, dla którego wszystkie drugie pochodne
zbiegają lokalnie jednostajnie do ciągłych funkcji.
Z podstawowego twierdzenia rachunku całkowego
na odcinkach równoległych do osi, przez przejście
do granicy najpierw dla $u_n$, a następnie dla
$\partial_i u_n$, wynika, że granice są odpowiednio
pierwszymi i drugimi pochodnymi klasycznymi $u$.
W szczególności $u\in C^2$. Nierówności Höldera
i oszacowania supremum przechodzą do granicy,
co daje deklarowaną normę $C^{2,\alpha}$.
Z równości $-\Delta u_n=f_n$ i lokalnie jednostajnej
zbieżności Hessianów wynika $-\Delta u=f$ klasycznie.
\end{proof}

\begin{twierdzenie}[Wewnętrzne oszacowanie Schaudera: model Poissona]
\label{thm:pde16-schauder-poisson}
Niech $0<\alpha<1$, $B_r\Subset B_R\subset\R^d$
będą współśrodkowymi kulami, a
$u\in C^2(B_R)\cap C^0(\overline{B_R})$ spełnia
$-\Delta u=f$ w $B_R$ dla $f\in C^{0,\alpha}(\overline{B_R})$.
Wówczas $u\in C^{2,\alpha}(\overline{B_r})$ i
\begin{equation}\label{eq:pde16-schauder-int}
 \|u\|_{C^{2,\alpha}(\overline{B_r})}
 \le C\bigl(\|u\|_{C^0(\overline{B_R})}
 +\|f\|_{C^{0,\alpha}(\overline{B_R})}\bigr).
\end{equation}
Stała zależy od $d,\alpha,r,R$.
\end{twierdzenie}
\begin{proof}
Przez translację zakładamy, że środek kul jest zerem.
Wybieramy $\rho=(r+R)/2$ i odcięcie
$\eta\in C_c^\infty(B_R)$ równe jeden na $B_\rho$,
z konstrukcji w lemacie~\ref{lem:pde16-localization}.
Niech $F$ będzie przedłużeniem $\eta f$ przez zero
na $\R^d$. Dla dwóch punktów w $B_R$ stosujemy
nierówność dla iloczynu z definicji
~\ref{def:pde16-holder}. Jeżeli $x\in B_R$ i
$y\notin B_R$, to $\eta(y)=0$ i
$|F(x)-F(y)|\le\|f\|_\infty|\eta(x)-\eta(y)|$.
Dla dwóch punktów poza $B_R$ różnica jest zerowa.
Gładkość i zwarty nośnik odcięcia dają skończoną
globalną półnormę $[\eta]_\alpha$ (dla odległości
co najwyżej jeden używamy gradientu, dla większych
normy supremum). Stąd
$\|F\|_{C^{0,\alpha}(\R^d)}
\le C_{d,\alpha,r,R}\|f\|_{C^{0,\alpha}(\overline{B_R})}$.

Z lematu~\ref{lem:pde16-newton-holder}
$v=\Gamma*F$ jest klasy $C^{2,\alpha}$ lokalnie,
$-\Delta v=F$ i pełna norma $C^{2,\alpha}$
na $\overline{B_\rho}$ jest kontrolowana przez
powyższą normę $F$.
W $B_\rho$ funkcja $h=u-v$ jest harmoniczna,
a z twierdzenia~\ref{thm:pde3-analytic} jest gładka.
Wszystkie jej pochodne są harmoniczne przez
przemienność pochodnych klasycznych. Ustalmy
$\sigma=(\rho-r)/4$ i zastosujmy oszacowanie
~\eqref{eq:pde3-derivative-bound}
do kul o promieniu $\sigma$ ze środkami
w $B_{\rho-\sigma}$. Otrzymujemy
$\|\partial_jh\|_{\infty,B_{\rho-\sigma}}
\le d\sigma^{-1}\|h\|_{\infty,B_\rho}$.
Następnie stosujemy to samo oszacowanie do pochodnych
pierwszego rzędu na $B_{\rho-2\sigma}$ i pochodnych
drugiego rzędu na $B_{\rho-3\sigma}$.
Ponieważ $r<\rho-3\sigma$, dla $1\le k\le3$ mamy
\[
 \|D^kh\|_{\infty,B_r}
 \le C_{d,k}(\rho-r)^{-k}\|h\|_{\infty,B_\rho}.
\]
Z wypukłości $B_r$ i całkowania wzdłuż odcinka
$[D^2h]_{\alpha,B_r}
\le (2r)^{1-\alpha}\|D^3h\|_{\infty,B_r}$
(z ewentualną stałą zależną od normy tensorowej).
Wreszcie
$\|h\|_{\infty,B_\rho}\le\|u\|_{\infty,B_R}
+\|v\|_{\infty,B_\rho}$.
Sumujemy oszacowania dla $u=h+v$,
włączając niższe pochodne i półnormę drugich pochodnych.
Ciągłość na domknięciu $B_r$ wynika z regularności
w jego otoczeniu. Otrzymujemy~\eqref{eq:pde16-schauder-int}
bez zakładania z góry regularności Hölderowskiej Hessianu.
\end{proof}

\begin{twierdzenie}[Skalarny współczynnik w postaci niedywergencyjnej]
\label{thm:pde16-scalar-schauder}
Niech $d\ge2$, $0<\alpha<1$, $\Omega\subset\R^d$ będzie obszarem,
$a,f\in C^{0,\alpha}_{\mathrm{loc}}(\Omega)$, przy czym
$a(x)\ge\lambda>0$ dla każdego $x\in\Omega$.
Jeżeli $u\in C^2(\Omega)$ spełnia $a(x)\Delta u=f$ w $\Omega$,
to $u\in C^{2,\alpha}_{\mathrm{loc}}(\Omega)$.
Dla niepustych otwartych zbiorów $U\Subset V\Subset\Omega$
zachodzi
\[
 \|u\|_{C^{2,\alpha}(\overline U)}
 \le C\bigl(\|f\|_{C^{0,\alpha}(\overline V)}
 +\|u\|_{C^0(\overline V)}\bigr),
\]
gdzie $C$ zależy tylko od $d,\alpha,U,V,\lambda$ i
$\|a\|_{C^{0,\alpha}(\overline V)}$.
\end{twierdzenie}
\begin{proof}
Na $\overline V$ funkcja $1/a$ należy do $C^{0,\alpha}$,
gdyż $\|1/a\|_\infty\le\lambda^{-1}$ oraz
$|a(x)^{-1}-a(y)^{-1}|
\le\lambda^{-2}[a]_\alpha|x-y|^\alpha$.
Z nierówności dla iloczynu wynika dla $g=f/a$
\[
 \|g\|_{C^{0,\alpha}(\overline V)}
 \le\lambda^{-1}\|f\|_{C^{0,\alpha}(\overline V)}
 +\lambda^{-2}\|f\|_\infty[a]_{C^{0,\alpha}(\overline V)}.
\]
Mamy $-\Delta u=-g$. Niech
$\ell=\operatorname{dist}(\overline U,\R^d\setminus V)/4>0$.
Dla każdego $x\in\overline U$ kula
$\overline{B_{2\ell}(x)}$ mieści się w $V$.
Twierdzenie~\ref{thm:pde16-schauder-poisson}
na parze $B_\ell(x),B_{2\ell}(x)$ daje jednolite
oszacowanie normy $C^{2,\alpha}$ na $B_\ell(x)$
przez $C(\|g\|_{C^{0,\alpha}(\overline V)}
+\|u\|_{\infty,V})$, niezależnie od $x$.
W szczególności wszystkie pochodne do rzędu dwa
są jednostajnie ograniczone na $\overline U$.
Dla $x,y\in\overline U$ z $0<|x-y|<\ell$
oba punkty leżą w $B_\ell(x)$, więc iloraz Höldera
pochodnych rzędu dwa ma powyższe ograniczenie.
Dla $|x-y|\ge\ell$ iloraz ograniczamy przez
$2\ell^{-\alpha}$ razy odpowiednią normę supremum.
To daje pełną normę na $\overline U$,
również gdy $U$ nie jest wypukły lub spójny.
Podstawiając oszacowanie $g$, otrzymujemy tezę.
Dowolność $U,V$ daje regularność lokalną.
\end{proof}

\begin{lemat}[Oszacowanie Schaudera \emph{a priori} i lokalizacja]
\label{lem:pde16-schauder-apriori}
Niech $0<\alpha<1$, $B_r\Subset B_R\Subset\Omega$
będą współśrodkowymi kulami, a
$A=(a_{ij})$, $b=(b_i)$, $c$ należą do
$C^{0,\alpha}(\overline{B_R})$; załóżmy, że $A$ jest
symetryczna i jednostajnie eliptyczna.
Jeżeli $v\in C^{2,\alpha}(\overline{B_R})$, to
\begin{equation}\label{eq:pde16-schauder-apriori}
 \|v\|_{C^{2,\alpha}(\overline{B_r})}
 \le C\bigl(\|A:D^2v+b\cdot\nabla v+cv\|_{C^{0,\alpha}(\overline{B_R})}
 +\|v\|_{C^0(\overline{B_R})}\bigr),
\end{equation}
gdzie $C$ zależy od $d,\alpha,r,R$, eliptyczności i norm
Hölderowskich współczynników, ale nie od $v$.
\end{lemat}
\begin{proof}
\emph{Krok 1: oszacowanie dla zamrożonego operatora.}
Dla kuli $B_\rho(x_0)$ i dowolnego pola $g$ oznaczamy
\[
 |g|_{\alpha,\rho;E}=\|g\|_{\infty,E}
                    +\rho^\alpha[g]_{\alpha,E},
 \qquad
 |w|_{k,\alpha,\rho;E}
 =\sum_{j=0}^k\rho^j\|D^jw\|_{\infty,E}
       +\rho^{k+\alpha}[D^kw]_{\alpha,E}.
\]
W tych oznaczeniach skala $\rho$ jest ustalona również
dla mniejszych zbiorów $E$. Dla iloczynów skalarnych
$|gh|_{\alpha,\rho;E}\le
|g|_{\alpha,\rho;E}|h|_{\alpha,\rho;E}$ na mocy
nierówności po definicji~\ref{def:pde16-holder}.
Dla iloczynów i kontrakcji pól o wartościach w skończenie
wymiarowych przestrzeniach wolno wprowadzić stałą
zależną jedynie od wymiaru.

Niech $w\in C^{2,\alpha}(\R^d)$ ma zwarty nośnik
w $B_\rho(x_0)$, a $A_0$ będzie stałą symetryczną
macierzą o wartościach własnych w $[\lambda,\Lambda]$.
Wówczas
\begin{equation}\label{eq:pde16-frozen}
 \rho^2|D^2w|_{\alpha,\rho;B_\rho(x_0)}
 \le C_0\rho^2|A_0:D^2w|_{\alpha,\rho;B_\rho(x_0)},
\end{equation}
gdzie $C_0$ zależy tylko od $d,\alpha,\lambda,\Lambda$.
Aby to wykazać, podstawiamy $x=x_0+\rho S z$,
$S=A_0^{1/2}$ i $\widetilde w(z)=w(x_0+\rho Sz)$.
Reguła łańcuchowa z kroku 1 dowodu twierdzenia
~\ref{thm:pde16-variable} daje
$\Delta_z\widetilde w=\rho^2(A_0:D^2w)(x_0+\rho Sz)$.
Nośnik $\widetilde w$ leży w kuli o promieniu
$\lambda^{-1/2}$, a
$B_\rho(x_0)$ po zmianie zmiennych mieści się
w tej samej kuli.
Dla $F=-\Delta_z\widetilde w$ mamy
$\Gamma*F=\widetilde w$: dla funkcji testowych wynika to
z lematu~\ref{lem:pde7-convolution-derivatives}
i $-\Delta\Gamma=\delta_0$, a dla $\widetilde w$
wygładzamy ją z jednostajną zbieżnością pochodnych
do rzędu dwa i wspólnym zwartym nośnikiem.
Lokalna całkowalność $\Gamma$ pozwala przejść
do granicy w całce potencjału.

Niech $G=A_0:D^2w$. Przedłużenie zerowe $G$
ma tę samą półnormę Höldera co ograniczenie do kuli:
dla pary punktów po przeciwnych stronach sfery
wybieramy punkt przecięcia łączącego je odcinka
ze sferą, gdzie $G=0$, i stosujemy oszacowanie
wewnątrz domkniętej kuli. Pary punktów poza kulą
dają zerową różnicę. Stąd
\[
 \|F\|_{C^{0,\alpha}(\R^d)}
 \le \rho^2\bigl(\|G\|_{\infty,B_\rho}
             +\Lambda^{\alpha/2}\rho^\alpha[G]_{\alpha,B_\rho}\bigr).
\]
Z lematu~\ref{lem:pde16-newton-holder} dla $F$
otrzymujemy oszacowanie normy $D_z^2\widetilde w$
w ustalonej kuli o promieniu $\lambda^{-1/2}+1$.
Z tożsamości
$D_x^2w(x_0+\rho Sz)=
\rho^{-2}S^{-\mathsf T}D_z^2\widetilde w(z)S^{-1}$
i nierówności
$|z-z'|\le\lambda^{-1/2}\rho^{-1}|x-x'|$
wynika~\eqref{eq:pde16-frozen}.

\emph{Krok 2: interpolacja na całej kuli.}
Wykażemy, że dla każdego $\varepsilon>0$ istnieje
$C_\varepsilon=C(d,\alpha,\varepsilon)$ takie, że
\[
 |v|_{\alpha,\rho;B_\rho}
 +\rho|\nabla v|_{\alpha,\rho;B_\rho}
 \le \varepsilon\rho^2\|D^2v\|_{\infty,B_\rho}
       +C_\varepsilon\|v\|_{\infty,B_\rho}.
\]
Wystarczy rozważyć kulę jednostkową, gdyż dla
$V(z)=v(x_0+\rho z)$ wszystkie składniki
przekształcają się zgodnie z powyższymi wagami.
Oznaczmy $Y=\|V\|_\infty$, $G=\|\nabla V\|_\infty$,
$H=\|D^2V\|_\infty$ na $B_1$.
Dla $x\in B_1$ i $0<t<1/2$ połóżmy $x_t=(1-t)x$.
Mamy $B_{t/2}(x_t)\subset B_1$.
Wzór Taylora z resztą całkową na odcinku
$x_t+[0,t/2]e_j$ daje
$|\partial_jV(x_t)|\le 4Y/t+(t/4)H$.
Ponadto z całkowania gradientu wzdłuż odcinka
$|\nabla V(x)-\nabla V(x_t)|\le tH$.
Po zsumowaniu składowych mamy
$G\le C_d(t^{-1}Y+tH)$.

Z wypukłości kuli
$|\nabla V(x)-\nabla V(y)|\le H|x-y|$,
a także ta różnica jest nie większa niż $2G$.
Dla liczb nieujemnych $a,b$ zachodzi
$\min(a,b)\le a^\alpha b^{1-\alpha}$.
Stąd $[\nabla V]_\alpha\le H^\alpha(2G)^{1-\alpha}$;
analogicznie $[V]_\alpha\le G^\alpha(2Y)^{1-\alpha}$.
W szczególności $[V]_\alpha\le G+2Y$.
Z ważonej nierówności Younga, dla każdego $\tau>0$,
mamy $H^\alpha(2G)^{1-\alpha}\le\tau H+C_{\alpha,\tau}G$.
Wybieramy najpierw $\tau=\varepsilon/2$, a następnie
$t$ tak małe, by po podstawieniu oszacowania $G$
całkowity współczynnik przy $H$ był nie większy
niż $\varepsilon$. Pozostałe wyrazy są ograniczone
przez $C_\varepsilon Y$. To dowodzi interpolacji
bez przedłużania funkcji poza kulę.

\emph{Krok 3: lokalne oszacowanie i pochłonięcie różnicy
współczynników.}
Piszemy $Lv=A:D^2v+b\cdot\nabla v+cv$ i wybieramy
kulę $\overline{B_\rho(x_0)}\subset B_R$, $\rho\le1$.
Niech $\eta$ będzie odcięciem równym jeden
na $B_{\rho/2}(x_0)$ i o zwartym nośniku w $B_\rho(x_0)$,
z konstrukcji w lemacie~\ref{lem:pde16-localization}.
Ta sama konstrukcja daje
$\|D^j\eta\|_\infty\le C_{d,j}\rho^{-j}$ dla $j\le3$.
Z wypukłości kuli i oszacowania pochodnej rzędu $j+1$
dla $j\le2$ wynika
$\rho^j|D^j\eta|_{\alpha,\rho;B_\rho}\le C_{d,\alpha}$.
Przedłużone przez zero $w=\eta v$ należy do
$C^{2,\alpha}(\R^d)$ i ma zwarty nośnik w $B_\rho$.
Istotnie, pochodne z reguły Leibniza znikają w otoczeniu
sfery, a półnormę dla par po przeciwnych stronach
sfery kontrolujemy przez punkt przecięcia odcinka
ze sferą jak w kroku 1.

Dla $A_0=A(x_0)$ mamy dokładną tożsamość
\[
 \begin{aligned}
 A_0:D^2w
 &=\eta Lv+(A_0-A):D^2w\\
 &\quad+2(A\nabla\eta)\cdot\nabla v+v A:D^2\eta\\
 &\quad-\eta b\cdot\nabla v-\eta cv .
 \end{aligned}
\]
Na $B_\rho$ zachodzi
$|A-A_0|_{\alpha,\rho;B_\rho}\le2\rho^\alpha[A]_{\alpha,B_R}$.
Zatem stosując~\eqref{eq:pde16-frozen} i nierówności
dla iloczynów, otrzymujemy dla $X=\rho^2|D^2w|_{\alpha,\rho;B_\rho}$
\[
 \begin{aligned}
 X &\le C_1\rho^2|Lv|_{\alpha,\rho;B_\rho}
       +C_1\rho^\alpha[A]_{\alpha,B_R} X\\
   &\quad+C_2\bigl(|v|_{\alpha,\rho;B_\rho}
                         +\rho|\nabla v|_{\alpha,\rho;B_\rho}\bigr).
 \end{aligned}
\]
Rzeczywiście, wyraz z $D\eta$ ma po przemnożeniu
przez $\rho^2$ czynnik co najwyżej $C\rho$,
wyraz z $D^2\eta$ czynnik $C$, a wyrazy z $b,c$
czynniki odpowiednio $C\rho^2$ przy normie gradientu
i $C\rho^2$ przy normie funkcji. Ponieważ $\rho\le1$,
ostatnie dwa wyrazy mieszczą się w zapisanym oszacowaniu.
Stałe $C_1,C_2$ zależą jedynie od danych twierdzenia,
nie od środka, promienia ani $v$.
Wybieramy $\rho_0\in(0,1]$ tak, by
$C_1\rho_0^\alpha[A]_{\alpha,B_R}\le1/2$.
Dla $\rho\le\rho_0$ pochłaniamy wyraz z $X$.
Na $B_{\rho/2}$ mamy $D^2w=D^2v$, a z kroku 2,
z parametrem odpowiednio zmniejszonym przez stałą $2C_2$,
otrzymujemy dla dowolnego $\varepsilon>0$
\[
 \begin{aligned}
 \rho^2|D^2v|_{\alpha,\rho;B_{\rho/2}}
 &\le C_3\rho^2|Lv|_{\alpha,\rho;B_\rho}
     +\varepsilon\rho^2\|D^2v\|_{\infty,B_\rho}\\
 &\quad+C_{4,\varepsilon}\|v\|_{\infty,B_\rho}.
 \end{aligned}
\]
Zauważmy, że mały składnik po prawej zawiera
jedynie normę supremum Hessianu, nie jego półnormę.

\emph{Krok 4: ważona norma i kontrola całego wnętrza.}
Dla $x\in B_R$ definiujemy
\[
 \delta(x)=\tfrac18\min\{\rho_0,
                 \operatorname{dist}(x,\partial B_R)\},\qquad
 M=\sup_{x\in B_R}\delta(x)^2|D^2v(x)|.
\]
Liczba $M$ jest skończona z założenia
$v\in C^{2,\alpha}(\overline{B_R})$.
Funkcja $\delta$ jest Lipschitzowska ze stałą $1/8$.
Kula $\overline{B_{2\delta(x)}(x)}$ mieści się w $B_R$,
jej promień nie przekracza $\rho_0$, a dla $y$
w tej kuli mamy $\delta(y)\ge3\delta(x)/4$.
Stąd
\[
 \delta(x)^2\|D^2v\|_{\infty,B_{2\delta(x)}(x)}
 \le (4/3)^2 M .
\]
W oszacowaniu z kroku 3 bierzemy $\rho=2\delta(x)$
i zachowujemy tylko normę supremum po lewej.
Po podzieleniu przez $4$ dostajemy
\[
 \delta(x)^2|D^2v(x)|
 \le C_\varepsilon(\|Lv\|_{C^{0,\alpha}(\overline{B_R})}
                  +\|v\|_{\infty,B_R})
       +\varepsilon(4/3)^2 M,
\]
ponieważ $\rho\le1$. Wybieramy na przykład
$\varepsilon=1/8$ i bierzemy supremum po $x$.
Współczynnik przy $M$ jest mniejszy niż $1/2$,
więc po pochłonięciu
\[
 M\le C(\|Lv\|_{C^{0,\alpha}(\overline{B_R})}
                  +\|v\|_{\infty,B_R})=:C Z.
\]
Na $\overline{B_r}$ zachodzi
$\delta(x)\ge\delta_*=\min\{\rho_0,R-r\}/8>0$.
W szczególności $\|D^2v\|_{\infty,B_r}\le C\delta_*^{-2}Z$.
Dla pary $x,y\in B_r$, $0<|x-y|<\delta_*$,
oba punkty leżą w $B_{\delta(x)}(x)$.
Ponowne zastosowanie kroku 3 z $\rho=2\delta(x)$
oraz oszacowanie ważonej normy dają
$\rho^{2+\alpha}[D^2v]_{\alpha,B_{\delta(x)}(x)}
\le C Z$. Iloraz Höldera tej pary jest więc
nie większy niż $C(2\delta_*)^{-2-\alpha}Z$.
Dla $|x-y|\ge\delta_*$ używamy dwukrotnej normy
supremum Hessianu pomnożonej przez $\delta_*^{-\alpha}$.
Otrzymujemy pełną kontrolę półnormy na $B_r$.
Z kroku 2 na $B_r$ wynika także oszacowanie gradientu
przez $\|D^2v\|_{\infty,B_r}$ i $\|v\|_{\infty,B_r}$,
ze stałą zależną od $r$. Przez ciągłość wszystkie
oszacowania obowiązują na domknięciu.
To dowodzi~\eqref{eq:pde16-schauder-apriori}.
\end{proof}

\begin{lemat}[Komutator wygładzania i mnożenia]
\label{lem:pde16-commutator}
Niech $0<\alpha<1$, $a\in C^{0,\alpha}(\R^d)$,
$g\in L^\infty(\R^d)$, a $\rho_\varepsilon(x)
=\varepsilon^{-d}\rho(x/\varepsilon)$, gdzie
$\rho\in C_c^\infty(B_1(0))$ jest nieujemna i ma całkę $1$.
Dla
\[
 \mathcal C_\varepsilon[a,g](x)
 =a(x)(g*\rho_\varepsilon)(x)
 -(ag)*\rho_\varepsilon(x)
\]
zachodzi jednostajnie względem $0<\varepsilon\le1$
\[
 \|\mathcal C_\varepsilon[a,g]\|_{C^{0,\alpha}(\R^d)}
 \le C_{d,\alpha,\rho}[a]_{C^{0,\alpha}(\R^d)}\|g\|_\infty.
\]
Ponadto $\|\mathcal C_\varepsilon[a,g]\|_\infty
\le \varepsilon^\alpha[a]_\alpha\|g\|_\infty$.
\end{lemat}
\begin{proof}
Oba sploty są dobrze określone w każdym punkcie:
całkujemy funkcje istotnie ograniczone z jądrem
z $L^1$. Zmiana reprezentanta $g$ na zbiorze miary zero
nie zmienia żadnej z tych całek.
Z definicji otrzymujemy
\[
 \mathcal C_\varepsilon[a,g](x)
 =\int_{\R^d}(a(x)-a(y))g(y)\rho_\varepsilon(x-y)
 \,\mathrm d\lambda_d(y).
\]
Na nośniku jądra $|x-y|\le\varepsilon$.
Z nieujemności jądra i jego całki równej jeden
wynika oszacowanie supremum
$\|\mathcal C_\varepsilon[a,g]\|_\infty
\le\varepsilon^\alpha[a]_\alpha\|g\|_\infty$.

Niech $h=|x-x'|>0$. Jeśli $h\ge\varepsilon$,
to różnica wartości komutatora jest co do modułu
nie większa niż
$2\varepsilon^\alpha[a]_\alpha\|g\|_\infty
\le2h^\alpha[a]_\alpha\|g\|_\infty$.
Dla $0<h<\varepsilon$ dodajemy i odejmujemy
$a(x')(g*\rho_\varepsilon)(x)$, uzyskując
\begin{align*}
 \mathcal C_\varepsilon[a,g](x)-\mathcal C_\varepsilon[a,g](x')
 &= (a(x)-a(x'))(g*\rho_\varepsilon)(x)\\
 &\quad+\int_{\R^d}(a(x')-a(y))g(y)
 [\rho_\varepsilon(x-y)-\rho_\varepsilon(x'-y)]
 \,\mathrm d\lambda_d(y).
\end{align*}
Pierwszy wyraz jest ograniczony przez
$[a]_\alpha\|g\|_\infty h^\alpha$.
W drugim $|x'-y|\le2\varepsilon$ wszędzie tam,
gdzie różnica jąder jest niezerowa.
Ponadto z podstawowego twierdzenia rachunku całkowego
wzdłuż odcinka między $x'-y$ i $x-y$ mamy
\[
 \begin{aligned}
 &\int_{\R^d}|\rho_\varepsilon(x-y)
               -\rho_\varepsilon(x'-y)|\,\mathrm d\lambda_d(y)\\
 &\quad\le h\int_0^1\int_{\R^d}
 |\nabla\rho_\varepsilon(x'+t(x-x')-y)|
 \,\mathrm d\lambda_d(y)\,\mathrm d\lambda_1(t)\\
 &\quad=h\|\nabla\rho_\varepsilon\|_1
 =h\varepsilon^{-1}\|\nabla\rho\|_1.
 \end{aligned}
\]
W drugiej linii zastosowano twierdzenie Tonellego,
a w trzeciej niezmienniczość miary Lebesgue'a
względem translacji oraz zmianę skali.
Zatem drugi wyraz jest ograniczony przez
$2^\alpha[a]_\alpha\|g\|_\infty\|\nabla\rho\|_1
\varepsilon^{\alpha-1}h$.
Ponieważ $\varepsilon^{\alpha-1}h
=h^\alpha(h/\varepsilon)^{1-\alpha}\le h^\alpha$,
otrzymujemy jednostajne oszacowanie półnormy Höldera.
Razem z normą supremum i $\varepsilon\le1$
daje ono deklarowaną pełną normę.
\end{proof}

\begin{twierdzenie}[Oszacowania Schaudera dla zmiennych współczynników]
\label{thm:pde16-schauder-variable}
Niech $0<\alpha<1$, $\Omega\subset\R^d$ będzie obszarem,
$a_{ij},b_i,c,f\in C^{0,\alpha}_{\mathrm{loc}}(\Omega)$,
a macierz $A=(a_{ij})$ będzie symetryczna i lokalnie
jednostajnie eliptyczna. Jeżeli $u\in C^2(\Omega)$ spełnia
\[
 Lu:=A:D^2u+b\cdot\nabla u+cu=f
 \qquad\text{w }\Omega,
\]
to $u\in C^{2,\alpha}_{\mathrm{loc}}(\Omega)$ oraz dla
niepustych otwartych zbiorów $U\Subset V\Subset\Omega$
\[
 \|u\|_{C^{2,\alpha}(\overline U)}
 \le C\bigl(\|f\|_{C^{0,\alpha}(\overline V)}
 +\|u\|_{C^0(\overline V)}\bigr),
\]
gdzie $C$ zależy od $U,V,d,\alpha$, norm
$C^{0,\alpha}(\overline V)$ współczynników i stałych
eliptyczności na $\overline V$, ale nie od $u,f$.
\end{twierdzenie}
\begin{proof}
\emph{Krok 1: lokalne przedłużenia i regularizacja.}
Wybierzmy współśrodkowe kule
$B_r\Subset B_R\Subset B_S\Subset\Omega$.
Niech $\chi\in C_c^\infty(B_S)$ będzie odcięciem
równym jeden na otoczeniu $\overline{B_R}$,
a $0\le\chi\le1$. Możemy użyć konstrukcji
z lematu~\ref{lem:pde16-localization}.
Niech $N$ będzie otwartym otoczeniem $\overline{B_R}$,
na którym $\chi=1$, zawartym w $B_S$.
Przedłużamy $\chi u$ przez zero na $\R^d$;
otrzymana funkcja $w$ jest klasy $C_c^2(\R^d)$,
gdyż $\chi$ znika w otoczeniu brzegu $B_S$.
Dla każdego współczynnika $a$ (dowolnej składowej $A,b,c$)
definiujemy $\widetilde a$ jako zerowe przedłużenie
$\chi a$, a $\widetilde f$ jako zerowe przedłużenie $\chi f$.
Funkcje te należą do $C_c^{0,\alpha}(\R^d)$:
wewnątrz $B_S$ stosujemy nierówność dla iloczynu
po definicji~\ref{def:pde16-holder}, natomiast
dla $x\in B_S$, $y\notin B_S$ mamy
$|\widetilde a(x)-\widetilde a(y)|
\le\|a\|_{\infty,B_S}|\chi(x)-\chi(y)|$.
Dla dwóch punktów poza $B_S$ różnica jest zerowa.
Analogicznie dla $\widetilde f$.
Lokalna regularność Höldera zapewnia skończoność
pełnych norm na $\overline{B_S}$: na zwartym
podzbiorze wybieramy skończone pokrycie kulami,
na których półnormy są skończone, i dla par punktów
bliskich stosujemy wspólne otoczenie z pokrycia,
a dla pozostałych ograniczenie supremum i dodatnią
liczbę Lebesgue'a pokrycia.

Kładziemy $u_\varepsilon=w*\rho_\varepsilon$.
Dla każdego wielowskaźnika $\beta$, $|\beta|\le2$,
$\partial^\beta u_\varepsilon
=(\partial^\beta w)*\rho_\varepsilon$:
różniczkowanie pod całką uzasadnia ograniczoność
ciągłych pochodnych $w$ i zwarty nośnik jądra.
Z ich jednostajnej ciągłości wynika
$\|\partial^\beta u_\varepsilon-\partial^\beta w\|_\infty
\to0$, ponieważ różnicę szacujemy przez supremum
$|\partial^\beta w(x-z)-\partial^\beta w(x)|$
dla $|z|\le\varepsilon$.
W szczególności $u_\varepsilon\to u$
w $C^2(\overline{B_R})$.

\emph{Krok 2: dokładne równanie po wygładzeniu.}
Bierzemy $0<\varepsilon<\min\{1,
\operatorname{dist}(\overline{B_R},\R^d\setminus N)\}$.
Dla $x\in\overline{B_R}$ nośnik jądra w zmiennej
całkowania $y$ leży w $N$, gdzie $w=u$ i
$\widetilde a=a$, $\widetilde f=f$.
W szczególności, dla $|\beta|\le2$,
\[
 a(x)\partial^\beta u_\varepsilon(x)
 =((\widetilde a\partial^\beta w)*\rho_\varepsilon)(x)
    +\mathcal C_\varepsilon[\widetilde a,\partial^\beta w](x).
\]
Sumując tożsamości dla wszystkich współczynników,
w pierwszych składnikach wykorzystujemy $Lu=f$
na $N$. Dostajemy na $\overline{B_R}$
\[
 \begin{aligned}
 Lu_\varepsilon
 &=\widetilde f*\rho_\varepsilon\\
 &\quad+\sum_{i,j=1}^d
       \mathcal C_\varepsilon[\widetilde a_{ij},\partial_{ij}w]\\
 &\quad+\sum_{i=1}^d
       \mathcal C_\varepsilon[\widetilde b_i,\partial_i w]
       +\mathcal C_\varepsilon[\widetilde c,w].
 \end{aligned}
\]
Współczynniki operatora po lewej stronie
nie zostały wygładzone.
Wszystkie drugie argumenty komutatorów są globalnie
ograniczone, gdyż $w\in C_c^2$.
Lemat~\ref{lem:pde16-commutator} daje jednostajne
oszacowania ich norm $C^{0,\alpha}(\R^d)$.
Także $\|\widetilde f*\rho_\varepsilon\|_\infty
\le\|\widetilde f\|_\infty$, a dla dowolnych $x,x'$
\[
 \begin{aligned}
 &|(\widetilde f*\rho_\varepsilon)(x)
       -(\widetilde f*\rho_\varepsilon)(x')|\\
 &\quad\le\int_{\R^d}
 |\widetilde f(x-z)-\widetilde f(x'-z)|
 \rho_\varepsilon(z)\,\mathrm d\lambda_d(z)
 \le[\widetilde f]_\alpha|x-x'|^\alpha .
 \end{aligned}
\]
Zatem prawa strona dokładnego równania jest
jednostajnie ograniczona w $C^{0,\alpha}(\overline{B_R})$.
Stała na tym etapie może zależeć od normy $C^2$
funkcji $w$, ale nie od $\varepsilon$.
Nie zakładamy Hölderowskiej ciągłości $D^2u$.

\emph{Krok 3: uzyskanie regularności.}
Każde $u_\varepsilon$ jest gładkie na $\R^d$,
a więc należy do $C^{2,\alpha}(\overline{B_R})$.
Możemy zastosować lemat
~\ref{lem:pde16-schauder-apriori} do $u_\varepsilon$,
z tymi samymi współczynnikami i stałą niezależną
od parametru wygładzania. Otrzymujemy
\[
 \|u_\varepsilon\|_{C^{2,\alpha}(\overline{B_r})}
 \le C\bigl(\|Lu_\varepsilon\|_{C^{0,\alpha}(\overline{B_R})}
                  +\|w\|_\infty\bigr)\le C_u.
\]
Dla każdych różnych $x,y\in\overline{B_r}$
przechodzimy do granicy w nierówności
$|D^2u_\varepsilon(x)-D^2u_\varepsilon(y)|
\le C_u|x-y|^\alpha$ dzięki jednostajnej zbieżności
drugich pochodnych. W ten sposób
$u\in C^{2,\alpha}(\overline{B_r})$.
Dowolność wewnętrznych kul daje
$u\in C^{2,\alpha}_{\mathrm{loc}}(\Omega)$.

\emph{Krok 4: oszacowanie niezależne od normy $C^2$ rozwiązania.}
W poprzednim kroku nie uzyskaliśmy jeszcze żądanej
zależności stałej, ponieważ oszacowanie komutatorów
zawierało normę $C^2$ funkcji $w$.
Teraz stosujemy lemat~\ref{lem:pde16-schauder-apriori}
bezpośrednio do $u$, którego regularność Höldera
już wykazano. Niech
$\ell=\operatorname{dist}(\overline U,\R^d\setminus V)/4$.
Dla każdego $x\in\overline U$ mamy
$\overline{B_{2\ell}(x)}\subset V$.
Na tej kuli $u$ należy do $C^{2,\alpha}$:
wynika to z uzyskanej regularności lokalnej
i zwartości domknięcia kuli, jak w kroku 1.
Ponieważ $Lu=f$, lemat daje
\[
 \|u\|_{C^{2,\alpha}(\overline{B_\ell(x)})}
 \le C\bigl(\|f\|_{C^{0,\alpha}(\overline V)}
                 +\|u\|_{\infty,V}\bigr).
\]
Stała jest niezależna od $x$: promienie są ustalone,
a normy współczynników i stałe eliptyczności
kontrolujemy na $\overline V$.
Otrzymujemy oszacowanie supremum wszystkich pochodnych
do rzędu dwa na $\overline U$.
Dla $x,y\in\overline U$, $0<|x-y|<\ell$, oba punkty
leżą w $B_\ell(x)$, gdzie obowiązuje powyższa
kontrola półnormy Höldera.
Dla $|x-y|\ge\ell$ iloraz Höldera ograniczamy
przez dwukrotne supremum odpowiedniej pochodnej
pomnożone przez $\ell^{-\alpha}$.
Sumujemy oszacowania składowych, otrzymując
żądaną pełną normę na $\overline U$.
\end{proof}

\begin{uwaga}[Granice formułowanych twierdzeń]
\label{rem:pde16-boundary}
W tym rozdziale udowodniliśmy oszacowania wewnętrzne,
nie zaś globalne twierdzenia regularności brzegowej
w $W^{2,p}$ i $C^{2,\alpha}$.
Typowymi założeniami w globalnej teorii Dirichleta
są brzeg klasy $C^{1,1}$ dla oszacowań $W^{2,p}$
oraz brzeg klasy $C^{2,\alpha}$ i dane brzegowe będące
śladem funkcji klasy $C^{2,\alpha}$ dla oszacowań
Schaudera; potrzebne są także odpowiednie założenia
o operatorze i prawej stronie aż do brzegu.
Nie są to twierdzenia o konieczności tych warunków
dla każdego rozwiązania.
Dla $p=2$ i postaci dywergencyjnej pełne twierdzenie
brzegowe z dokładnymi hipotezami udowodniono
w rozdziale poprzednim (twierdzenie~\ref{thm:pde15-boundary}).
Dowód ogólnych globalnych oszacowań wymaga odrębnej
kontroli przy brzegu, której nie zastępuje rezultat
wewnętrzny: stałe tego ostatniego mogą rosnąć
przy zmniejszaniu odległości od brzegu.
\end{uwaga}

\begin{uwaga}[Postać równania i przestrzeń rozwiązania]
\label{rem:pde-v42-two-regularity-classes}
Twierdzenie~\ref{thm:pde15-boundary} dotyczy słabych
rozwiązań $u\in H_0^1(\Omega)$ równania w postaci
dywergencyjnej $-\diver(A\nabla u)=f$ na ograniczonym
obszarze klasy $C^{1,1}$, przy jednostajnej eliptyczności,
$A\in W^{1,\infty}$ i $f\in L^2$.
Natomiast twierdzenie~\ref{thm:pde16-schauder-variable}
dotyczy rozwiązań klasy $C^2$ równania
w postaci niedywergencyjnej $Lu=f$ przy lokalnie
jednostajnie eliptycznej symetrycznej macierzy $A$
oraz lokalnie Hölderowskich współczynnikach i prawej
stronie, a jego wniosek jest lokalny.
Nie wolno przenosić tezy $C^{2,\alpha}$ na słabe
rozwiązania postaci dywergencyjnej z jedynie
$C^{0,\alpha}$-współczynnikami: rozwinięcie dywergencji
do równania niedywergencyjnego wymaga osobnego
uzasadnienia, gdyż pochodne współczynników nie muszą
być funkcjami ograniczonymi. Samo oszacowanie
$W^{2,p}$ z twierdzenia~\ref{thm:pde16-variable}
również nie zastępuje argumentu uzyskania regularności.
\end{uwaga}

\section{Kontrprzykłady i granice regularności}
\begin{kontrprzyklad}[Brak ciągłości $D^2u$ przy ograniczonej prawej stronie]
Na przedziale $(-1,1)$ niech $u(x)=-\tfrac12x|x|$.
Dla $x>0$ mamy $u(x)=-x^2/2$, a dla $x<0$
mamy $u(x)=x^2/2$. W zerze iloraz różnicowy
$u(h)/h=-|h|/2$ dąży do zera.
Zatem $u$ jest klasy $C^1$, a $u'(x)=-|x|$ na całym
przedziale. Funkcja $u'$ jest Lipschitzowska i jej
słaba pochodna wynosi $-\operatorname{sgn}(x)$:
dla testu $\varphi\in C_c^\infty((-1,1))$
całkujemy przez części na obu półprzedziałach,
a wyrazy w zerze znikają, ponieważ $u'(0)=0$.
W konsekwencji
\[
 -u''=\operatorname{sgn}(x)\quad\text{prawie wszędzie},
 \qquad u\in W^{2,\infty}((-1,1)).
\]
Na ograniczonym przedziale wynika stąd również
$u\in W^{2,p}((-1,1))$ dla każdego $1\le p<\infty$.
Jednostronne ilorazy różnicowe funkcji $u'$ w zerze
mają granice $-1$ i $1$, więc druga pochodna klasyczna
w zerze nie istnieje. Tym samym $u\notin C^2((-1,1))$.
W wymiarze $d\ge2$ na kostce $(-1,1)^d$ bierzemy
$U(x)=-x_1|x_1|/2$. Z twierdzenia Fubiniego
i jednowymiarowego rachunku słabych pochodnych
wynika $\partial_{11}U=-\operatorname{sgn}(x_1)$
oraz zerowe pozostałe drugie pochodne.
Otrzymujemy $-\Delta U=\operatorname{sgn}(x_1)$,
$U\in W^{2,\infty}$, ale $U\notin C^2$.
\end{kontrprzyklad}

\begin{kontrprzyklad}[Ostrość wykładnika Höldera]
Dla $0<\alpha<1$ połóżmy na $(-1,1)$
$u(x)=-|x|^{2+\alpha}/((1+\alpha)(2+\alpha))$.
Dla $x\ne0$ bezpośrednio otrzymujemy
\[
 u'(x)=-\frac{\operatorname{sgn}(x)|x|^{1+\alpha}}{1+\alpha},
 \qquad u''(x)=-|x|^\alpha.
\]
Iloraz $u(h)/h$ dąży do zera, więc $u'(0)=0$.
Następnie $u'(h)/h=-|h|^\alpha/(1+\alpha)\to0$,
co daje $u''(0)=0$ i ciągłość drugiej pochodnej.
Zatem $-\Delta u=|x|^\alpha$ (w wymiarze jeden)
także w sensie klasycznym w zerze.

Dla $s,t\ge0$ zachodzi $(s+t)^\alpha\le s^\alpha+t^\alpha$:
jeśli $s>0$, wynika to z monotoniczności funkcji
$z\mapsto (1+z)^\alpha-z^\alpha$ dla $z>0$,
której pochodna jest ujemna, a granica przy $z\downarrow0$
wynosi $1$; przypadek $s=0$ daje równość.
Stąd dla $x,y\in[-1,1]$,
\[
 \bigl||x|^\alpha-|y|^\alpha\bigr|
 \le\bigl||x|-|y|\bigr|^\alpha
 \le |x-y|^\alpha.
\]
Zatem $u\in C^{2,\alpha}([-1,1])$.
Dla $\alpha<\beta\le1$ mamy natomiast
$|u''(h)-u''(0)|/h^\beta=h^{\alpha-\beta}\to+\infty$
przy $h\downarrow0$, więc druga pochodna nie jest
Hölderowska z wykładnikiem $\beta$ nawet lokalnie
w otoczeniu zera. W wymiarze $d\ge2$ funkcja
$U(x)=u(x_1)$ na kostce $(-1,1)^d$ daje ten sam
kontrprzykład dla laplasjanu: jedyną niezerową
składową Hessianu jest $u''(x_1)$, a ilorazy
wzdłuż pierwszej osi mają powyższą postać.
\end{kontrprzyklad}

\section{Podsumowanie i miejsce w dalszej teorii}
Twierdzenie~\ref{thm:pde16-cz} kontroluje drugie
pochodne potencjału Newtona przez normę prawej
strony w $L^p$ dla $1<p<\infty$; nie zapewnia
całkowalności samego potencjału na $\R^d$.
Silne oszacowania dla $p=1$ i $p=\infty$ zawodzą
w ogólności (uwaga~\ref{rem:pde16-endpoints}),
a oszacowanie~\eqref{eq:pde16-weak11} dotyczy słabego
typu $(1,1)$ na $L^1\cap L^2$.
Decydujące są znoszenie się jądra o średniej zero,
rozkład Calderóna--Zygmunda, interpolacja i dualność.
Twierdzenie~\ref{thm:pde16-variable} dla ciągłych
współczynników jest oszacowaniem a priori
w klasie $W^{2,p}_{\mathrm{loc}}$, nie samodzielnym
twierdzeniem o istnieniu lub uzyskaniu tej regularności.

W teorii Schaudera odejmowanie $f(x)$ pod całką
osobliwą daje zbieżność dzięki Hölderowskiemu
oszacowaniu różnicy. Twierdzenia
~\ref{thm:pde16-schauder-poisson}
i~\ref{thm:pde16-schauder-variable} dają wewnętrzną
regularność $C^{2,\alpha}$ dla $0<\alpha<1$
rozwiązań klasy $C^2$ przy odpowiednich hipotezach.
Nie obejmują końca $\alpha=1$ ani dowolnych
ograniczonych prawych stron. Nie dają też
oszacowania globalnego aż do brzegu.
W kolejnym rozdziale~\ref{chap:pde17} rozpatrzymy
regularność De Giorgiego--Nasha--Mosera dla słabych
rozwiązań skalarnych jednorodnych równań eliptycznych
w postaci dywergencyjnej, przy ograniczonych mierzalnych
i jednostajnie eliptycznych współczynnikach.
Regularność Hölderowska dotyczy tam samego rozwiązania,
nie automatycznie jego drugich pochodnych.
